\documentclass{article}

\usepackage{arxiv}

\usepackage[utf8]{inputenc} % allow utf-8 input
\usepackage[T1]{fontenc}    % use 8-bit T1 fonts
\usepackage{amsmath}        % math environments (cases, dfrac, binom, text, boldsymbol)
\usepackage{amssymb}        % extended math symbols
\usepackage{amsfonts}       % blackboard math symbols
\usepackage{hyperref}       % hyperlinks
\usepackage{url}            % simple URL typesetting
\usepackage{booktabs}       % professional-quality tables
\usepackage{nicefrac}       % compact symbols for 1/2, etc.
\usepackage{microtype}      % microtypography
\usepackage{cleveref}       % smart cross-referencing
\usepackage{lipsum}         % Can be removed after putting your text content
\usepackage{graphicx}
\usepackage{natbib}
\usepackage{doi}
\usepackage{placeins}

\title{Structure and Dynamics of the Urban Interpersonal Attitude Network}

% Here you can change the date presented in the paper title
%\date{September 9, 1985}
% Or remove it
%\date{}

\newif\ifuniqueAffiliation
% Comment to use multiple affiliations variant of author block 
\uniqueAffiliationtrue

\ifuniqueAffiliation % Standard variant of author block
\author{ \href{https://orcid.org/0000-0003-0417-8108}{\IfFileExists{orcid.pdf}{\includegraphics[scale=0.06]{orcid.pdf}\hspace{1mm}}{}César Rennó-Costa}\thanks{Use footnote for providing further
		information about author (webpage, alternative
		address)---\emph{not} for acknowledging funding agencies.} \\
	Digital Metropolis Institute (IMD)\\
	Federal University of Rio Grande do Norte (UFRN)\\
	Natal, RN, Brazil\\
	\texttt{cesar@imd.ufrn.br} \\
	%% examples of more authors
	%\And
	%\href{https://orcid.org/0000-0000-0000-0000}{\includegraphics[scale=0.06]{orcid.pdf}\hspace{1mm}Elias D.~Striatum} \\
	%Department of Electrical Engineering\\
	%Mount-Sheikh University\\
	%Santa Narimana, Levand \\
	%\texttt{stariate@ee.mount-sheikh.edu} \\
	%% \AND
	%% Coauthor \\
	%% Affiliation \\
	%% Address \\
	%% \texttt{email} \\
	%% \And
	%% Coauthor \\
	%% Affiliation \\
	%% Address \\
	%% \texttt{email} \\
	%% \And
	%% Coauthor \\
	%% Affiliation \\
	%% Address \\
	%% \texttt{email} \\
}
\else
% Multiple affiliations variant of author block
\usepackage{authblk}
\renewcommand\Authfont{\bfseries}
\setlength{\affilsep}{0em}
% box is needed for correct spacing with authblk
\newbox{\orcid}\sbox{\orcid}{\includegraphics[scale=0.06]{orcid.pdf}} 
\author[1]{%
	\href{https://orcid.org/0000-0000-0000-0000}{\usebox{\orcid}\hspace{1mm}David S.~Hippocampus\thanks{\texttt{hippo@cs.cranberry-lemon.edu}}}%
}
\author[1,2]{%
	\href{https://orcid.org/0000-0000-0000-0000}{\usebox{\orcid}\hspace{1mm}Elias D.~Striatum\thanks{\texttt{stariate@ee.mount-sheikh.edu}}}%
}
\affil[1]{Department of Computer Science, Cranberry-Lemon University, Pittsburgh, PA 15213}
\affil[2]{Department of Electrical Engineering, Mount-Sheikh University, Santa Narimana, Levand}
\fi

% Uncomment to override  the `A preprint' in the header
%\renewcommand{\headeright}{Technical Report}
%\renewcommand{\undertitle}{Technical Report}
\renewcommand{\shorttitle}{Structure and Dynamics of the Urban Interpersonal Attitude Network}

%%% Add PDF metadata to help others organize their library
%%% Once the PDF is generated, you can check the metadata with
%%% $ pdfinfo template.pdf
\hypersetup{
pdftitle={Structure and Dynamics of the Urban Interpersonal Attitude Network},
pdfsubject={physics.soc-ph, stat.AP, q-bio.PE},
pdfauthor={César Rennó-Costa},
pdfkeywords={attitude dynamics, multiplex social networks, ideological alignment, opinion formation, urban scaling},
}

\begin{document}
\maketitle

\begin{abstract}
	\lipsum[1]
\end{abstract}


% keywords can be removed
\keywords{First keyword \and Second keyword \and More}


\section{Introduction}
\lipsum[2]
\lipsum[3]




% ==============================================================================
% Methods Section: Multidimensional Population Structure, Interpersonal
% Multiplex Network, Mass Media, and Algorithmic Platforms
% ==============================================================================

\section{Methods}
\label{sec:methods_baseline}

To establish a comprehensive mathematical and computational framework governing public attitude formation, we formulate a discrete-time, multidimensional agent-based model of an urban population interacting concurrently across three distinct communication substrates: a five-layer physical and social multiplex network, an institutional traditional mass media broadcast ecosystem, and an algorithmic social media platform network. We allocate all state arrays in contiguous double-precision memory and integrate them synchronously using parallelized, just-in-time (JIT)-compiled numerical kernels \citep{lam2015numba}.

% ------------------------------------------------------------------------------
\subsection{Overall Simulation Architecture and Multidimensional Concept Space}
\label{subsec:concept_space}

We consider a city populated by $N$ citizens indexed by $i \in \mathcal{V} = \{1, 2, \dots, N\}$, evolving over discrete time steps $t \in \{0, 1, \dots, T\}$ with integration step size $\Delta t = 0.15$. To ensure statistical robustness and exact replicability across high-performance computing (HPC) array runs, we evaluate each experimental condition across $S = 100$ independent Monte Carlo realizations ($s \in \{0, \dots, 99\}$). For each realization $s$, an independent cryptographic seed sequence spawns dedicated pseudorandom number generator (PRNG) streams for demographic initialization, network topology construction, and stochastic interpersonal dynamics.

Public issues and policy concepts are embedded in a compact, $D$-dimensional continuous attitude space
\begin{equation}
    \Omega = [-1, 1]^D, \qquad D = 4,
    \label{eq:concept_space}
\end{equation}
where each coordinate $d \in \{0, 1, 2, 3\}$ represents a primary policy domain. Coordinate $d = 0$ captures economic and fiscal policy (spanning state-led redistribution at $-1$ to market deregulation at $+1$), $d = 1$ spans moral and social customs (progressive secular norms at $-1$ to traditionalist conservatism at $+1$), $d = 2$ represents environment and climate (preservation regulation at $-1$ to resource extraction at $+1$), and $d = 3$ models public security and institutional order (civil liberties at $-1$ to punitive law enforcement at $+1$). The origin $\mathbf{0} = (0, 0, 0, 0)^\top \in \Omega$ represents a neutral, centrist position across all four thematic dimensions. For any two vectors $\mathbf{u}, \mathbf{v} \in \Omega$, we define the dimension-normalized Euclidean metric
\begin{equation}
    d(\mathbf{u}, \mathbf{v}) = \frac{1}{\sqrt{D}} \|\mathbf{u} - \mathbf{v}\|_2 = \sqrt{\frac{1}{D} \sum_{d=0}^{D-1} (u_d - v_d)^2} \in [0, 2],
    \label{eq:normalized_metric}
\end{equation}
so that the distance of any vector $\mathbf{u} \in \Omega$ from the ideological center $\mathbf{0}$ lies strictly in the unit interval
\begin{equation}
    E(\mathbf{u}) \equiv d(\mathbf{u}, \mathbf{0}) = \frac{1}{\sqrt{D}} \|\mathbf{u}\|_2 \in [0, 1].
    \label{eq:distance_to_center}
\end{equation}

Because citizens do not evaluate policy domains in complete isolation, logical and ideological interdependence across the $D = 4$ dimensions is represented by a row-normalized cognitive coupling matrix $\mathbf{C} = [C_{dd'}] \in \mathbb{R}^{D \times D}$ \citep{converse1964nature, friedkin2016network}. Specifically, $\mathbf{C}$ combines dominant diagonal self-consistency with decaying off-diagonal cross-issue constraints
\begin{equation}
    \tilde{C}_{dd'} =
    \begin{cases}
        1, & \text{if } d = d', \\[4pt]
        (-1)^{d + d'} \dfrac{0.18}{|d - d'|}, & \text{if } d \neq d',
    \end{cases}
    \qquad
    C_{dd'} = \frac{\tilde{C}_{dd'}}{\max\!\left(1, \sum_{k=0}^{D-1} |\tilde{C}_{dk}|\right)}.
    \label{eq:coupling_matrix}
\end{equation}

% ------------------------------------------------------------------------------
\subsection{Citizen Demographics and Behavioral Characteristics}
\label{subsec:demographics_traits}

Each citizen $i \in \mathcal{V}$ has static demographic attributes calibrated to the Brazilian adult population ($\ge 16$ years) from the 2022 Demographic Census \citep{ibge2022censo}, along with heterogeneous psychological traits that govern attitude plasticity.

\paragraph{Age Distribution ($a_i$).}
Citizen age $a_i \in [16, 90]$ is sampled from the six contiguous adult age cohorts of the 2022 IBGE Census pyramid: $B_1 = [16, 25)$ ($p_1 = 0.168$), $B_2 = [25, 35)$ ($p_2 = 0.198$), $B_3 = [35, 45)$ ($p_3 = 0.196$), $B_4 = [45, 55)$ ($p_4 = 0.164$), $B_5 = [55, 65)$ ($p_5 = 0.138$), and $B_6 = [65, 90]$ ($p_6 = 0.136$). Within cohorts $B_1, \dots, B_5$, raw ages are drawn uniformly within bracket bounds, whereas within the senior cohort $B_6$ raw ages follow a mortality-tapered Beta distribution $a_i^{\mathrm{raw}} = 65 + 25\,\mathrm{Beta}(1.15, 2.25)$. To avoid artificial discretization discontinuities at cohort boundaries, we apply a zero-mean Gaussian smoothing kernel $\xi_i \sim \mathcal{N}(0, \sigma_a^2)$ with $\sigma_a = 1.35 \, \mathrm{yr}$, reflected at the domain boundaries $\{16, 90\}$. Each citizen's normalized life-stage coordinate is defined as
\begin{equation}
    \tilde{a}_i = \frac{a_i - 16}{90 - 16} = \frac{a_i - 16}{74} \in [0, 1].
    \label{eq:age_norm}
\end{equation}
For cohort-level analysis, citizens are grouped into Young ($a_i < 30$), Mid-Age ($30 \le a_i < 55$), and Senior ($a_i \ge 55$) life stages.

\paragraph{Socioeconomic Status ($\mathrm{SES}_i$), Gender ($g_i$), and Residential District ($r_i$).}
Socioeconomic and educational status $\mathrm{SES}_i \in [0, 1]$ is drawn independently from a right-skewed Beta distribution
\begin{equation}
    \mathrm{SES}_i \sim \mathrm{Beta}(2.0, 3.0), \qquad \mathbb{E}[\mathrm{SES}_i] = 0.40,
    \label{eq:ses_distribution}
\end{equation}
reflecting the empirical skew of urban income and educational attainment in Brazil \citep{ibge2022censo}. Citizens are stratified into Low SES ($\mathrm{SES}_i < 0.35$, $\approx 42.8\%$), Mid SES ($0.35 \le \mathrm{SES}_i < 0.55$, $\approx 35.1\%$), and High SES ($\mathrm{SES}_i \ge 0.55$, $\approx 22.1\%$). Gender $g_i \in \{0, 1, 2\}$ is sampled with probabilities $\{0.48, 0.50, 0.02\}$, and each citizen is assigned to one of $R = 10$ urban residential districts $r_i \in \{0, \dots, R-1\}$ according to city-specific Dirichlet weights $\mathbf{\pi}_r \sim \mathrm{Dirichlet}(3\mathbf{1}_R)$.

\paragraph{Convictions ($\mathbf{s}_i(t)$) and Current Opinions ($\mathbf{x}_i(t)$).}
Following belief-system theory \citep{converse1964nature, zaller1992nature, kiley2020measuring}, each citizen $i$ maintains two coupled attitudinal vectors in $\Omega = [-1, 1]^D$, comprising a slow-moving deep conviction $\mathbf{s}_i(t) = (s_{i,0}(t), \dots, s_{i,D-1}(t))^\top$ that anchors personal worldview and crystallizes over time \citep{friedkin1990social, alwin1991aging}, and a fast-moving current opinion $\mathbf{x}_i(t) = (x_{i,0}(t), \dots, x_{i,D-1}(t))^\top$ that responds to daily interpersonal deliberation and external events while remaining elastically tethered to $\mathbf{s}_i(t)$ \citep{zaller1992nature}. The individual's conviction extremity at time $t$ is defined by the normalized distance of their conviction from the origin \citep{howe2017attitude}
\begin{equation}
    E_i(t) \equiv E(\mathbf{s}_i(t)) = \frac{1}{\sqrt{D}} \|\mathbf{s}_i(t)\|_2 \in [0, 1].
    \label{eq:conviction_extremity}
\end{equation}
Citizens are classified by conviction extremity into Centrists ($E_i < 0.35$), Moderates ($0.35 \le E_i \le 0.60$), and Extremists ($E_i > 0.60$).

\paragraph{Individual Psychological Parameters ($\alpha_i, \epsilon_i, u_i, b_i$).}
Guided by the aging-stability hypothesis \citep{alwin1991aging, kiley2020measuring} and attitude-strength theory \citep{sherif1961social, howe2017attitude}, four behavioral traits govern citizen plasticity. Social susceptibility $\alpha_i \in [0.05, 0.98]$ dictates the relative weight placed on interpersonal social influence versus internal conviction anchoring \citep{friedkin1990social},
\begin{equation}
    \alpha_i = \mathrm{clip}_{[0.05,\, 0.98]}\!\left(0.85 - 0.55\tilde{a}_i - 0.22 E_i(0) + \zeta_{i,\alpha}\right),
    \label{eq:susceptibility}
\end{equation}
where $\zeta_{i,\alpha} \sim \mathcal{U}(-0.08, 0.08)$ (or $0.85 - 0.65\tilde{a}_i + \zeta_{i,\alpha}$ under unconstrained uniform initialization). The bounded confidence threshold $\epsilon_i \in [0.10, 1.50]$ (baseline $\bar{\epsilon} = 0.45$) defines the maximum attitudinal distance within which citizen $i$ accepts peer influence \citep{deffuant2000mixing, hegselmann2002opinion},
\begin{equation}
    \epsilon_i = \mathrm{clip}_{[0.12,\, 1.20]}\!\left(\bar{\epsilon}\,[1.15 - 0.45 E_i(0)] + \zeta_{i,\epsilon}\right).
    \label{eq:bounded_confidence}
\end{equation}
Attitude uncertainty $u_i(t) \in [0.02, 1.0]$ modulates relative agreement in peer encounters \citep{deffuant2002how}, initialized as
\begin{equation}
    u_i(0) = \mathrm{clip}_{[0.05,\, 0.95]}\!\left(0.65 - 0.35 E_i(0) - 0.20\tilde{a}_i + \zeta_{i,u}\right).
    \label{eq:initial_uncertainty}
\end{equation}
Finally, biased assimilation $b_i \in [0.0, 5.0]$ (baseline $\bar{b} = 1.25$) quantifies confirmatory evaluation toward directionally aligned peers \citep{lord1979biased, dandekar2013biased},
\begin{equation}
    b_i = \mathrm{clip}_{[0.10,\, 4.00]}\!\left(\bar{b}\,[0.65 + 0.75 E_i(0)] + \zeta_{i,b}\right).
    \label{eq:assimilation_bias}
\end{equation}

\paragraph{Initial Attitude Distributions.}
To evaluate convergence and stability from contrasting structural baselines, three initial condition regimes are defined at $t = 0$ with $\mathbf{x}_i(0) = \mathbf{s}_i(0)$. Under the uniform random baseline (\texttt{uniform\_random}), components satisfy $s_{i,d}(0) \stackrel{\text{i.i.d.}}{\sim} \mathcal{U}(-1, 1)$ across all citizens and dimensions, providing an unconstrained maximum-entropy null model ($\mathrm{Var}[s_{i,d}(0)] = 1/3$, $\mathbb{E}[\|\mathbf{s}_i(0)\|_2/\sqrt{D}] \approx 0.577$). Under the centrist baseline (\texttt{centrist\_convictions}), components follow $s_{i,d}(0) \sim \mathrm{clip}_{[-1,1]}(\mathcal{N}(0, 0.25^2))$, representing a predominantly moderate society. Under the extremist baseline (\texttt{extremist\_convictions}), $82\%$ of citizens are initialized around opposing ideological poles $p_i (-1)^d (0.72) + \mathcal{N}(0, 0.16^2)$ with $p_i \in \{-1, +1\}$, while $18\%$ occupy the ideological center $\mathcal{N}(0, 0.16^2)$.

% ------------------------------------------------------------------------------
\subsection{Real-World Five-Layer Multiplex Interpersonal Network}
\label{subsec:multiplex_network}

Citizens interact through a five-layer physical and social multiplex network $\mathcal{G} = (\mathcal{V}, \{\mathcal{E}^{[\ell]}\}_{\ell=1}^5)$ \citep{boccaletti2014structure, kivela2014multilayer}, where each layer $\ell \in \{1, \dots, 5\}$ represents a distinct sphere of everyday urban interaction with unweighted adjacency matrix $\mathbf{A}^{[\ell]} \in \{0, 1\}^{N \times N}$ and relative daily interaction weight $w_\ell > 0$ ($\sum_{\ell=1}^5 w_\ell = 1$).

In Layer 1 (Household Cliques, $\ell = 1$, $w_1 = 0.36$), the population $\mathcal{V}$ is partitioned into disjoint, fully connected household cliques of size $s_h \in \{2, 3, 4, 5\}$ drawn with probabilities $\{0.42, 0.35, 0.16, 0.07\}$, calibrated to the 2022 IBGE Census mean of $2.79$ residents per household \citep{ibge2022censo}. Because household assignment cuts across age brackets, this layer provides high-trust cross-generational ties. In Layer 2 (Workplace and School Peers, $\ell = 2$, $w_2 = 0.26$, $\langle k^{[2]} \rangle = 10$), citizens are ordered by a composite demographic key $\vartheta_i = 2.5 \lfloor(a_i - 16)/12\rfloor + \mathrm{SES}_i + \mathcal{N}(0, 0.08^2)$ and linked to their $\langle k^{[2]} \rangle / 2 = 5$ nearest peers on a circular lattice, capturing age-stratified classrooms for youth aged $16$ to $24$ \citep{inep2023censo} and workplace establishments for adults aged $25$ to $64$ \citep{mte2023rais}. In Layer 3 (Religious Congregation and Local Community, $\ell = 3$, $w_3 = 0.14$, $\langle k^{[3]} \rangle = 8$), citizens are sorted by their moral conviction coordinate perturbed by mixing noise, $\tilde{s}_{i,1} = s_{i,1}(0) + \mathcal{N}(0, 0.15^2)$, and linked to $\langle k^{[3]} \rangle = 8$ neighbors, producing community clusters with moderate ideological homophily \citep{mcpherson2001birds} that cut across generations. In Layer 4 (Small-World Friendship Circle, $\ell = 4$, $w_4 = 0.14$, $\langle k^{[4]} \rangle = 6$), ties are generated via the Watts--Strogatz mechanism \citep{watts1998collective} over a ring ordered by the multidimensional projection $\mathbf{s}_i(0)^\top \mathbf{\nu}$, with base coordination number $k_{\mathrm{ws}} = 6$ and long-range rewiring probability $p_{\mathrm{ws}} = 0.15$. Finally, Layer 5 (Transit Weak Ties, $\ell = 5$, $w_5 = 0.10$, $\langle k^{[5]} \rangle = 4$) is constructed as an Erd\H{o}s--R\'enyi random graph with $N \langle k^{[5]} \rangle / 2$ undirected edges connecting citizens uniformly at random across districts, capturing weak cross-cutting encounters in shared public transit \citep{granovetter1973strength, ibge2022censo}.

The five layers are aggregated into a single row-normalized composite interpersonal adjacency matrix $\mathbf{W} = [W_{ij}] \in [0, 1]^{N \times N}$ stored in Compressed Sparse Row (CSR) format
\begin{equation}
    \tilde{W}_{ij} = \sum_{\ell=1}^5 w_\ell A_{ij}^{[\ell]}, \qquad
    W_{ij} = \frac{\tilde{W}_{ij}}{\sum_{k \in \mathcal{V}} \tilde{W}_{ik}},
    \label{eq:composite_csr}
\end{equation}
yielding a mean composite interpersonal degree of $\langle k_{\mathrm{comp}} \rangle \approx 30.8$ connections per citizen. Demographic mixing across the six IBGE age cohorts $B_1, \dots, B_6$ is measured by the $6 \times 6$ Age-Cohort Connection Matrix $\mathbf{M} = [M_{ab}]$
\begin{equation}
    M_{ab} = \frac{\sum_{i \in B_a} \sum_{j \in B_b} \mathbb{I}(W_{ij} > 0)}{\sum_{i \in B_a} \sum_{k \in \mathcal{V}} \mathbb{I}(W_{ik} > 0)} \times 100\%,
    \label{eq:age_cohort_matrix}
\end{equation}
which records the percentage of composite ties originating from cohort $a$ that connect to cohort $b$.

% ------------------------------------------------------------------------------
\subsection{Individual Behavior: Coupled Conviction and Opinion Dynamics}
\label{subsec:dynamics_equations}

At each discrete step $t \to t + 1$, citizen attitudes evolve through synchronous interpersonal social influence over the composite network $\mathbf{W}$, coupled with exogenous fact assimilation, elastic conviction tethering, and slow conviction crystallization.

\subsubsection{Interpersonal Social Influence on Current Opinions}
For each citizen $i \in \mathcal{V}$ and each connected peer $j \in \mathcal{N}_i = \{j : W_{ij} > 0\}$, the normalized attitudinal distance between current opinions is
\begin{equation}
    d_{ij}(t) = \frac{1}{\sqrt{D}} \|\mathbf{x}_j(t) - \mathbf{x}_i(t)\|_2.
    \label{eq:peer_dist}
\end{equation}
Peer $j$ influences citizen $i$ at step $t$ if and only if $d_{ij}(t)$ falls within citizen $i$'s uncertainty-modulated bounded confidence threshold \citep{deffuant2000mixing, deffuant2002how, hegselmann2002opinion}
\begin{equation}
    d_{ij}(t) \le \tilde{\epsilon}_i(t), \qquad \tilde{\epsilon}_i(t) = \epsilon_i + 0.5\,u_i(t).
    \label{eq:effective_threshold}
\end{equation}
When Eq.~\eqref{eq:effective_threshold} is satisfied, interaction weight integrates relative agreement overlap, biased assimilation, and topological tie strength. Relative agreement overlap \citep{deffuant2002how} ensures peers with lower uncertainty $u_j(t)$ exert stronger influence proportional to uncertainty overlap $h_{ij}(t)$
\begin{equation}
    h_{ij}(t) = \min\!\big(\tilde{\epsilon}_i(t), u_j(t)\big) - \max\!\big(-\tilde{\epsilon}_i(t), d_{ij}(t) - u_j(t)\big), \qquad
    \phi_{ij}^{\mathrm{RA}}(t) = \max\!\left(0, \frac{h_{ij}(t)}{2u_j(t) + 10^{-9}}\right).
    \label{eq:relative_agreement}
\end{equation}
Biased assimilation \citep{lord1979biased, dandekar2013biased} provides an exponential confirmatory boost for directionally aligned peers according to cosine similarity $\cos\theta_{ij}(t)$
\begin{equation}
    \cos\theta_{ij}(t) = \frac{\mathbf{x}_i(t)^\top \mathbf{x}_j(t)}{\|\mathbf{x}_i(t)\|_2 \|\mathbf{x}_j(t)\|_2 + 10^{-12}}, \qquad
    \phi_{ij}^{\mathrm{BA}}(t) = \frac{\exp\!\big(0.35\,b_i \cos\theta_{ij}(t)\big)}{1 + 0.35\,b_i}.
    \label{eq:biased_assimilation}
\end{equation}
Combining structural weight $W_{ij}$ with relative agreement and biased assimilation yields the effective peer weight
\begin{equation}
    \tilde{w}_{ij}(t) = W_{ij} \big(0.4 + 0.6\,\phi_{ij}^{\mathrm{RA}}(t)\big) \phi_{ij}^{\mathrm{BA}}(t) \, \mathbb{I}\!\big(d_{ij}(t) \le \tilde{\epsilon}_i(t)\big).
    \label{eq:effective_peer_weight}
\end{equation}

The normalized interpersonal social shift vector $\Delta \mathbf{x}_i^{\mathrm{soc}}(t) \in \mathbb{R}^D$ is computed as the weighted average of peer differences, transformed by cognitive coupling $\mathbf{C}$ (Eq.~\eqref{eq:coupling_matrix})
\begin{equation}
    \Delta \mathbf{x}_i^{\mathrm{soc}}(t) =
    \begin{cases}
        \displaystyle \frac{\sum_{j \in \mathcal{N}_i} \tilde{w}_{ij}(t) \big(\mathbf{x}_j(t) - \mathbf{x}_i(t)\big)}{\sum_{j \in \mathcal{N}_i} \tilde{w}_{ij}(t)}, & \text{if } \sum_{j \in \mathcal{N}_i} \tilde{w}_{ij}(t) > 0, \\[10pt]
        \mathbf{0}, & \text{otherwise}.
    \end{cases}
    \label{eq:social_shift}
\end{equation}
Following the Friedkin--Johnsen framework \citep{friedkin1990social, parsegov2017novel}, the interpersonal update balances social pull $\mathbf{C}\Delta \mathbf{x}_i^{\mathrm{soc}}(t)$ against conviction anchoring $\mathbf{s}_i(t)$
\begin{equation}
    \mathbf{x}_i^{\mathrm{int}}(t) = \mathrm{clip}_{[-1, 1]^D}\!\Big(
        \mathbf{x}_i(t) + \Delta t \big[(1 - \alpha_i)(\mathbf{s}_i(t) - \mathbf{x}_i(t)) + \alpha_i \mathbf{C}\Delta \mathbf{x}_i^{\mathrm{soc}}(t)\big] + \boldsymbol{\xi}_i(t)
    \Big),
    \label{eq:fj_interpersonal_update}
\end{equation}
where $\boldsymbol{\xi}_i(t) \sim \mathcal{N}(\mathbf{0}, \sigma_{\xi}^2 \mathbf{I}_D)$ with $\sigma_{\xi} = 0.012$ represents idiosyncratic daily attitude noise. Concurrently, attitude uncertainty $u_i(t)$ adapts toward accepted peer confidence \citep{deffuant2002how}
\begin{equation}
    u_i(t+1) = \mathrm{clip}_{[0.02,\, 1.0]}\!\left(
        u_i(t) + 0.25\,\Delta t \frac{\sum_{j \in \mathcal{N}_i} \tilde{w}_{ij}(t)\big(u_j(t) - u_i(t)\big)}{\sum_{j \in \mathcal{N}_i} \tilde{w}_{ij}(t)}
    \right).
    \label{eq:uncertainty_update}
\end{equation}

\subsubsection{Exogenous Facts, Elastic Conviction Tethering, and Slow Conviction Crystallization}
\label{subsubsec:facts_and_tether}

\paragraph{Definition of Exogenous Facts.}
Beyond peer deliberation, citizens encounter objective real-world conditions, termed \textbf{exogenous facts} \citep{mullainathan2005market, gentzkow2006media}. At step $t$, available facts are represented by points $\mathbf{f}_m(t) \in \Omega = [-1, 1]^D$ ($m = 1, \dots, M_t$) with salience weights $\omega_m > 0$. Under stationary conditions (\texttt{neutral} regime), facts are generated symmetrically around the origin, $\mathbf{f}_m(t) \sim \mathrm{clip}_{[-0.85, 0.85]^D}(\mathcal{N}(\mathbf{b}_{\mathrm{fact}}(t), \sigma_f^2 \mathbf{I}_D))$ with $\mathbf{b}_{\mathrm{fact}}(t) = \mathbf{0}$ and $\sigma_f = 0.32$. Cross-cutting interpersonal mixing and unbiased facts naturally draw initially dispersed attitudes toward a centrist distribution around $\mathbf{0}$.

When the exogenous fact environment displays a directional impulse $\mathbf{b}_{\mathrm{fact}}(t) \in [-1, 1]^D$ indicated by $I_{\mathrm{fact}}(t) \in \{0, 1\}$, citizen $i$'s current opinion responds rapidly with sensitivity $\gamma_i(t)$ modulated by conviction extremity $E_i(t)$, age $\tilde{a}_i$, and socioeconomic status $\mathrm{SES}_i$ \citep{zaller1992nature, howe2017attitude}
\begin{equation}
    \gamma_i(t) = \gamma_0 \big(1.30 - 0.72 E_i(t)\big)\big(1.18 - 0.48\tilde{a}_i\big)\big(0.85 + 0.30\,\mathrm{SES}_i\big) I_{\mathrm{fact}}(t),
    \label{eq:fact_direct_pull}
\end{equation}
\begin{equation}
    \tilde{\mathbf{x}}_i(t) = \mathbf{x}_i^{\mathrm{int}}(t) + \gamma_i(t)\big(\mathbf{b}_{\mathrm{fact}}(t) - \mathbf{x}_i^{\mathrm{int}}(t)\big),
    \label{eq:opinion_fact_response}
\end{equation}
where $\gamma_0 = 0.055$ is baseline fact pull. Equation~\eqref{eq:fact_direct_pull} formalizes three empirical regularities: centrist citizens ($E_i \ll 1$) respond more strongly to new evidence than ideological extremists ($E_i \to 1$) \citep{howe2017attitude}; younger adults exhibit higher plasticity than older cohorts \citep{alwin1991aging}; and higher-SES citizens maintain higher direct informational exposure \citep{zaller1992nature}.

\paragraph{Elastic Conviction Tether.}
While current opinion $\tilde{\mathbf{x}}_i(t)$ responds rapidly to daily encounters and facts, it cannot drift arbitrarily far from deep conviction $\mathbf{s}_i(t)$ within short timeframes \citep{sherif1961social, zaller1992nature}. For raw excursion distance $g_i(t) = \frac{1}{\sqrt{D}}\|\tilde{\mathbf{x}}_i(t) - \mathbf{s}_i(t)\|_2$, each citizen has a maximum opinion excursion radius $\Delta_{\max, i}(t)$ around $\mathbf{s}_i(t)$
\begin{equation}
    \Delta_{\max, i}(t) = \max\!\Big(0.05,\; \Delta_0 \big(0.65 + 0.55\alpha_i\big) \big[\max(0.18, 1.15 - 0.95 E_i(t))\big]^{1.35} \big(1.16 - 0.32\tilde{a}_i\big)\Big),
    \label{eq:tether_radius}
\end{equation}
where $\Delta_0 = 0.34$ is base excursion radius. The opinion at step $t+1$ is smoothly bounded within radius $\Delta_{\max, i}(t)$ via a radial hyperbolic-tangent tether
\begin{equation}
    g_i^{\mathrm{eff}}(t) = \Delta_{\max, i}(t) \tanh\!\left(\frac{g_i(t)}{\Delta_{\max, i}(t)}\right), \qquad
    \mathbf{x}_i(t+1) = \mathrm{clip}_{[-1, 1]^D}\!\left(\mathbf{s}_i(t) + \frac{g_i^{\mathrm{eff}}(t)}{g_i(t) + 10^{-12}}\big(\tilde{\mathbf{x}}_i(t) - \mathbf{s}_i(t)\big)\right).
    \label{eq:elastic_tether}
\end{equation}
For small perturbations ($g_i(t) \ll \Delta_{\max, i}(t)$), $\mathbf{x}_i(t+1) \approx \tilde{\mathbf{x}}_i(t)$ moves freely; as $g_i(t)$ approaches $\Delta_{\max, i}(t)$, the tether smoothly saturates excursion at $\Delta_{\max, i}(t)$.

\paragraph{Slow Conviction Adaptation and Crystallization.}
Deep convictions $\mathbf{s}_i(t)$ evolve on a slower timescale than current opinions and progressively crystallize over time \citep{alwin1991aging, kiley2020measuring}. The conviction adaptation rate $\eta_i(t)$ combines a decaying crystallization term $\eta_{\mathrm{cryst}, i}(t)$ governed by characteristic timescale $\tau_c = 280$ steps and a nonlinear sustained-tension term $\eta_{\mathrm{sust}, i}(t)$ active when current opinion is held persistently near the tether boundary
\begin{equation}
    \kappa(t) = \frac{1}{1 + (t / \tau_c)^{1.15}},
    \label{eq:crystallization_factor}
\end{equation}
\begin{equation}
    \eta_{\mathrm{cryst}, i}(t) = \eta_0 \,\kappa(t) \big(1.25 - 0.75\tilde{a}_i\big)\big(1.22 - 0.68 E_i(t)\big)\big(1.10 - 0.25\,\mathrm{SES}_i\big),
    \label{eq:eta_cryst}
\end{equation}
\begin{equation}
    \eta_{\mathrm{sust}, i}(t) = 0.28\,\eta_0 \left(\frac{g_i^{\mathrm{eff}}(t)}{\Delta_{\max, i}(t)}\right)^2 \big(1.18 - 0.58\tilde{a}_i\big)\big(1.20 - 0.62 E_i(t)\big),
    \label{eq:eta_sustained}
\end{equation}
\begin{equation}
    \mathbf{s}_i(t+1) = \mathrm{clip}_{[-1, 1]^D}\!\Big(\mathbf{s}_i(t) + \big[\eta_{\mathrm{cryst}, i}(t) + \eta_{\mathrm{sust}, i}(t)\big]\big(\mathbf{x}_i(t+1) - \mathbf{s}_i(t)\big)\Big),
    \label{eq:conviction_update}
\end{equation}
where $\eta_0 = 0.0024 \ll \Delta t$ is baseline conviction drift rate. Because $\eta_0$ is two orders of magnitude smaller than interpersonal rates, transient opinion fluctuations are filtered out, while prolonged shifts slowly pull underlying convictions over hundreds of steps.

\begin{table}[htbp]
\centering
\caption{Summary of population demographics, interpersonal multiplex network, and two-timescale conviction--opinion parameters.}
\label{tab:model_parameters}
\footnotesize
\setlength{\tabcolsep}{4.5pt}
\begin{tabular}{lp{5.1cm}p{3.6cm}p{3.8cm}}
\toprule
\textbf{Symbol} & \textbf{Description} & \textbf{Value / Range} & \textbf{Reference} \\
\midrule
\multicolumn{4}{l}{\textit{Global Simulation \& Concept Space}} \\
$N$ & City population size & $\{10^4, 10^5, 10^6, 10^7\}$ & \citet{bettencourt2007growth} \\
$D$ & Number of policy/concept dimensions & $4$ ($\Omega = [-1, 1]^4$) & \citet{converse1964nature} \\
$T, \Delta t$ & Horizon and integration step size & $T \in \{100, 1500\}$, $\Delta t = 0.15$ & \citet{friedkin1990social} \\
$S$ & Monte Carlo realizations per scenario & $100$ & This study \\
$\mathbf{C}$ & Cognitive issue-coupling matrix & $\mathbb{R}^{4 \times 4}$, Eq.~\eqref{eq:coupling_matrix} & \citet{friedkin2016network} \\
\midrule
\multicolumn{4}{l}{\textit{Citizen Demographics \& Psychological Traits}} \\
$a_i, \tilde{a}_i$ & Citizen age (years) and normalized age & $[16, 90]$, $\tilde{a}_i \in [0, 1]$ & \citet{ibge2022censo} \\
$\mathrm{SES}_i$ & Socioeconomic and educational index & $\mathrm{Beta}(2, 3) \in [0, 1]$ & \citet{ibge2022censo} \\
$\mathbf{s}_i(t)$ & Deep conviction vector of citizen $i$ & $[-1, 1]^D$ & \citet{converse1964nature, kiley2020measuring} \\
$\mathbf{x}_i(t)$ & Current opinion vector of citizen $i$ & $[-1, 1]^D$ & \citet{zaller1992nature} \\
$E_i(t)$ & Conviction extremity ($\|\mathbf{s}_i(t)\|_2 / \sqrt{D}$) & $[0, 1]$ & \citet{howe2017attitude} \\
$\alpha_i$ & Interpersonal social susceptibility & $[0.05, 0.98]$, Eq.~\eqref{eq:susceptibility} & \citet{friedkin1990social, alwin1991aging} \\
$\epsilon_i$ & Bounded confidence (base $\bar{\epsilon}=0.45$) & $[0.10, 1.50]$, Eq.~\eqref{eq:bounded_confidence} & \citet{deffuant2000mixing, hegselmann2002opinion} \\
$u_i(t)$ & Attitude uncertainty (Relative Agreement) & $[0.02, 1.00]$, Eq.~\eqref{eq:uncertainty_update} & \citet{deffuant2002how} \\
$b_i$ & Biased assimilation (base $\bar{b}=1.25$) & $[0.00, 5.00]$, Eq.~\eqref{eq:assimilation_bias} & \citet{lord1979biased, dandekar2013biased} \\
\midrule
\multicolumn{4}{l}{\textit{Five-Layer Multiplex Interpersonal Network}} \\
$w_1, \langle s_h \rangle$ & Layer 1 (\texttt{family}): household clique & $w_1 = 0.36$, $\langle s_h \rangle = 2.79$ & \citet{ibge2022censo} \\
$w_2, \langle k^{[2]} \rangle$ & Layer 2 (\texttt{work\_school}): peer degree & $w_2 = 0.26$, $\langle k^{[2]} \rangle = 10$ & \citet{inep2023censo, mte2023rais} \\
$w_3, \langle k^{[3]} \rangle$ & Layer 3 (\texttt{church\_community}): degree & $w_3 = 0.14$, $\langle k^{[3]} \rangle = 8$ & \citet{mcpherson2001birds} \\
$w_4, \langle k^{[4]} \rangle, p_{\mathrm{ws}}$ & Layer 4 (\texttt{friends}): small-world circle & $w_4 = 0.14$, $\langle k^{[4]} \rangle = 6$, $p_{\mathrm{ws}}=0.15$ & \citet{watts1998collective} \\
$w_5, \langle k^{[5]} \rangle$ & Layer 5 (\texttt{weak\_ties}): transit ties & $w_5 = 0.10$, $\langle k^{[5]} \rangle = 4$ & \citet{granovetter1973strength, ibge2022censo} \\
\midrule
\multicolumn{4}{l}{\textit{Two-Timescale Conviction--Opinion \& Fact Parameters}} \\
$\Delta_0, \Delta_{\max, i}$ & Base and individual opinion tether radius & $\Delta_0 = 0.34$, Eq.~\eqref{eq:tether_radius} & \citet{sherif1961social, howe2017attitude} \\
$\eta_0$ & Base conviction adaptation drift rate & $0.0024 \text{ step}^{-1}$ & \citet{kiley2020measuring} \\
$\tau_c$ & Conviction crystallization timescale & $280 \text{ steps}$ & \citet{alwin1991aging} \\
$\gamma_0$ & Direct exogenous fact pull coefficient & $0.055 \text{ step}^{-1}$ & \citet{mullainathan2005market, zaller1992nature} \\
\bottomrule
\end{tabular}
\end{table}

% ------------------------------------------------------------------------------
\subsection{Traditional Mass Media Ecosystem and Exposure Dynamics}
\label{subsec:methods_traditional_media}

To capture the broadcast synchronization of public attention and institutional agenda setting, we introduce a competitive traditional media ecosystem comprising $M$ media conglomerates \citep{mullainathan2005market, gentzkow2006media, mccombs1972agenda}. Each conglomerate $g \in \{1, \dots, M\}$ operates an editorial portfolio across five distinct dissemination modalities $v \in \{0, 1, 2, 3, 4\}$ calibrated to empirical communications registers in Brazil \citep{kantar2023inside, reuters2023digital}.

Broadcast Television ($v = 0$) provides authoritative linear news broadcasts with high institutional credibility ($\gamma_0 = 0.75$) and intermediate novelty decay ($\delta_0 = 0.30$), reaching an empirical audience heavily concentrated in older age cohorts,
\begin{equation}
    P_{\mathrm{tv}}(a_i, \mathrm{SES}_i) = \mathrm{clip}_{[0.05,\, 0.85]}\!\left(0.24 + 0.48\,\tilde{a}_i^{0.9} - 0.10\,\mathrm{SES}_i\right).
    \label{eq:reach_tv}
\end{equation}
Radio AM/FM ($v = 1$) delivers rapid daily news bulletins with intermediate credibility ($\gamma_1 = 0.65$) and higher news-cycle turnover ($\delta_1 = 0.45$), targeted at working-age commuters and middle socioeconomic strata,
\begin{equation}
    P_{\mathrm{radio}}(a_i, \mathrm{SES}_i) = \mathrm{clip}_{[0.05,\, 0.75]}\!\left(0.18 + 0.34\,\tilde{a}_i + 0.14\sin(\pi\,\mathrm{SES}_i)\right).
    \label{eq:reach_radio}
\end{equation}
Written Press ($v = 2$) comprises investigative print journalism and established news portals with maximal institutional credibility ($\gamma_2 = 0.82$) and slow archival decay ($\delta_2 = 0.18$), strongly concentrated among higher-SES citizens and college graduates \citep{reuters2023digital},
\begin{equation}
    P_{\mathrm{press}}(a_i, \mathrm{SES}_i) = \mathrm{clip}_{[0.01,\, 0.65]}\!\left(0.02 + 0.42\,\mathrm{SES}_i^{1.8} + 0.12\,\tilde{a}_i\right).
    \label{eq:reach_press}
\end{equation}
Outdoor Marketing ($v = 3$) captures out-of-home billboards and transit signage functioning as an ambient urban presence with minimal attention decay ($\delta_3 = 0.08$) and moderate credibility ($\gamma_3 = 0.45$), reaching mobile commuting adults across transit corridors,
\begin{equation}
    P_{\mathrm{outdoor}}(a_i, \mathrm{SES}_i) = \mathrm{clip}_{[0.10,\, 0.70]}\!\left(0.35 + 0.25\left(1 - \frac{|\tilde{a}_i - 0.45|}{0.45}\right) + 0.05\,\mathrm{SES}_i\right).
    \label{eq:reach_outdoor}
\end{equation}
Finally, Public Gatherings and Civic Rallies ($v = 4$) encompass episodic political rallies and municipal events characterized by intense emotional contagion ($\gamma_4 = 0.80$) but rapid dissipation ($\delta_4 = 0.70$), engaging younger and politically active cohorts,
\begin{equation}
    P_{\mathrm{rally}}(a_i, \mathrm{SES}_i) = \mathrm{clip}_{[0.02,\, 0.30]}\!\left(0.04 + 0.14\,(1 - \tilde{a}_i) + 0.08\,(1 - \mathrm{SES}_i)\right).
    \label{eq:reach_rally}
\end{equation}

\paragraph{Editorial Centers, Tendentiousness, and Notice Projection.}
Each conglomerate $g$ maintains an institutional ideological center of attitude $\mathbf{m}_g \in \Omega$ and an editorial tendentiousness parameter $\tau_g \in [0, 1]$ \citep{mullainathan2005market, gentzkow2006media}. When an exogenous factual event $\mathbf{f}_k(t) \in \Omega$ arrives at time $t$, outlet $v$ belonging to group $g$ projects the raw fact into a published notice $\mathbf{y}_{k,g,v}(t) \in \Omega$,
\begin{equation}
    \mathbf{y}_{k,g,v}(t) = (1 - \eta_{g,v})\,\mathbf{f}_k(t) + \eta_{g,v}\,\mathbf{m}_g + \boldsymbol{\xi}_{k,g,v}, \qquad \eta_{g,v} = \mathrm{clip}_{[0, 1]}(\tau_g \mu_v),
    \label{eq:notice_projection_methods}
\end{equation}
where $\mu_v \in (0, 1]$ is the vehicle's baseline dissemination reach and $\boldsymbol{\xi}_{k,g,v} \sim \mathcal{N}(\mathbf{0}, \sigma_{\mathrm{media}}^2 \mathbf{I}_D)$ ($\sigma_{\mathrm{media}} = 0.03$) represents journalistic reporting noise. Under factual reporting ($\tau_g = 0.10$), published notices cluster tightly around the objective event $\mathbf{f}_k(t)$. Under tendentious reporting ($\tau_g = 0.90$), gatekeeping collapses coverage onto the outlet's ideological center $\mathbf{m}_g$ \citep{puglisi2015newspaper}. Each notice carries an active broadcast intensity $I_v(t)$ that decays exponentially according to its modality-specific rate,
\begin{equation}
    I_v(t+1) = I_v(t)\,\exp(-\delta_v).
    \label{eq:media_intensity_decay}
\end{equation}

\paragraph{Selective Exposure and Cognitive Media Assimilation.}
Conditional on accessing modality $v$ via Eqs.~\eqref{eq:reach_tv}--\eqref{eq:reach_rally}, citizen $i$ selects an active vehicle within that modality via distance-weighted softmax selectivity,
\begin{equation}
    q_{i,v}(t) = \frac{I_v(t)\,\big(d(\mathbf{x}_i(t), \mathbf{y}_v(t)) + \delta_{\mathrm{reg}}\big)^{-\beta_m}}{\sum_{v' \in \mathcal{V}_m} I_{v'}(t)\,\big(d(\mathbf{x}_i(t), \mathbf{y}_{v'}(t)) + \delta_{\mathrm{reg}}\big)^{-\beta_m}},
    \label{eq:vehicle_selection}
\end{equation}
where $\beta_m = 1.2$ for television, radio, and press, $\beta_{\mathrm{outdoor}} = 0.15$ for unselective ambient exposure, and $\beta_{\mathrm{rally}} = 2.16$ for partisan rally sorting. Exposure induces attitudinal change if the notice falls within citizen $i$'s credibility-expanded bounded confidence bound,
\begin{equation}
    d(\mathbf{x}_i(t), \mathbf{y}_v(t)) \le \epsilon_i^{\mathrm{eff}}(t) \equiv \epsilon_i \big(1 + 1.10\,\gamma_v\big),
    \label{eq:media_bounded_conf}
\end{equation}
yielding an opinion update towards the vehicle's report,
\begin{equation}
    \Delta \mathbf{x}_{i,v}^{\mathrm{trad}}(t) = \kappa_0 \gamma_v \alpha_i \left(0.35 + 0.65\left(1 - \frac{d(\mathbf{x}_i, \mathbf{y}_v)}{\epsilon_i^{\mathrm{eff}}}\right)\right) \big(1 - e^{-I_v(t)}\big) \exp\!\big(0.25\,b_i (0.5 - d(\mathbf{x}_i, \mathbf{y}_v))\big) \big(\mathbf{y}_v(t) - \mathbf{x}_i(t)\big),
    \label{eq:traditional_opinion_shift}
\end{equation}
where $\kappa_0 = 0.05$ is baseline broadcast strength. Concurrently, repeated institutional media exposure cultivates and internalizes into underlying convictions \citep{mccombs1972agenda, dellavigna2007fox},
\begin{equation}
    \mathbf{s}_i(t+1) = \mathrm{clip}_{[-1, 1]^D}\!\Big(\mathbf{s}_i(t) + 0.35\,\mathbf{C}\Delta \mathbf{x}_{i,v}^{\mathrm{trad}}(t)\Big),
    \label{eq:media_conviction_cultivation}
\end{equation}
while slightly reducing attitudinal uncertainty, $u_i(t+1) = \max(0.05, u_i(t)[1 - 0.008])$.

% ------------------------------------------------------------------------------
\subsection{Algorithmic Social Media Platforms, Influencer Hierarchies, and Cascades}
\label{subsec:methods_social_media}

To capture digital communication spaces, our framework integrates four distinct social media platform modalities alongside physical multiplex networks and mass media \citep{cinelli2021echo, tucker2018social}.

Visual Influencer Feeds (Instagram and TikTok, $p = 0$) feature algorithmic recommendation reels and creator feeds characterized by rapid content discovery, exhibiting steep youth concentration \citep{kantar2023inside, reuters2023digital},
\begin{equation}
    P_{\mathrm{IG/TT}}(a_i, \mathrm{SES}_i) = \mathrm{clip}_{[0.08,\, 0.94]}\!\left(0.84 - 0.62\,\tilde{a}_i^{0.85} + 0.10\,\mathrm{SES}_i\right).
    \label{eq:reach_ig_tt}
\end{equation}
Video Commentary and Microblogging (YouTube and X, $p = 1$) encompass analytical channels and public microblogging engaging working-age and politically attentive cohorts,
\begin{equation}
    P_{\mathrm{YT/X}}(a_i, \mathrm{SES}_i) = \mathrm{clip}_{[0.12,\, 0.88]}\!\left(0.56 - 0.28\,\tilde{a}_i + 0.24\,\mathrm{SES}_i^{1.2}\right).
    \label{eq:reach_yt_x}
\end{equation}
Social Graph Feeds (Facebook, $p = 2$) provide algorithmic timeline feeds anchored in peer network ties, peaking among mature and older adults,
\begin{equation}
    P_{\mathrm{FB}}(a_i, \mathrm{SES}_i) = \mathrm{clip}_{[0.10,\, 0.78]}\!\left(0.22 + 0.44\sin(\pi\,\tilde{a}_i^{0.75}) - 0.08\,\mathrm{SES}_i\right).
    \label{eq:reach_fb}
\end{equation}
Closed Messaging Networks (WhatsApp and Telegram, $p = 3$) represent high-penetration ($P_{\mathrm{WA}} \ge 0.80$), end-to-end encrypted private groups of bounded capacity ($K \le 256$) supporting viral multi-hop dark-social forwarding \citep{resende2019whatsapp, garimella2020images}.

\paragraph{Emergence of Scale-Free Influencer Hierarchies.}
A subset of digital users act as specialized content creators (influencers) indexed by $k \in \{1, \dots, N_{\mathrm{inf}}\}$. Influencer follower accumulation follows a competitive fitness-mediated preferential attachment process \citep{barabasi1999emergence, bianconi2001competition},
\begin{equation}
    \Pi_k \propto k_i \cdot \eta_k,
    \label{eq:preferential_attachment_methods}
\end{equation}
where $\eta_k = (0.65 + 0.70\,\tau_{\mathrm{inf}, k})$ represents creator fitness modulated by ideological tendentiousness $\tau_{\mathrm{inf}, k} \in [0.10, 0.90]$. The stationary follower in-degree distribution converges to a heavy-tailed power law,
\begin{equation}
    P(k) \sim k^{-\alpha}, \qquad \alpha \approx 1.15, \quad \mathrm{Gini} \approx 0.73,
    \label{eq:follower_power_law}
\end{equation}
wherein the top $5\%$ of elite influencers command $54.5\%$ of aggregate audience reach \citep{bakshy2011everyone}. In addition to organic broadcast reach, creator networks undergo coevolution through two coupled mechanisms. Influencers gain followers proportional to current reach $k_i$, creator fitness $\eta_k$, and algorithmic viral signals $R_m$,
\begin{equation}
    \Delta k_k^+(t) = \left(1 - \frac{k_k(t)}{K_{\mathrm{cap}}}\right) \Big[0.032\ln(1 + k_k)\,\eta_k \ln(1 + R_m) + 0.055\,\mathrm{Likes}_k\Big].
    \label{eq:follower_gain}
\end{equation}
Concurrently, dissonance-induced unfollowing occurs when attitudinal distance between a follower and creator exceeds tolerance, causing followers to sever ties at rate $p_{\mathrm{unfollow}} \propto d(\mathbf{x}_i, \mathbf{x}_k) - \epsilon_i$ and accelerating ideological sorting into digital echo chambers \citep{sasahara2021social}.

\paragraph{Closed-Loop Algorithmic Recommendation and Dual-Valence Engagement.}
Recommendation algorithms select content items $j$ for user $i$ across active platforms by optimizing for user retention and content engagement \citep{brady2017emotion, vosoughi2018spread},
\begin{equation}
    q_{ij}(t) \propto \big(d(\mathbf{x}_i(t), \mathbf{y}_j(t)) + \delta\big)^{-\beta} (1 + \mu_{\mathrm{emo}} e_j) \ln(1 + \mathrm{pop}_j) \big(0.5 + 0.5\ln(1 + k_{\mathrm{inf}, j})\big),
    \label{eq:recom_prob}
\end{equation}
where $\beta = 2.0$ governs homophilous filtering, $e_j \in [0, 1]$ is sensationalism and emotional charge ($\mu_{\mathrm{emo}} = 1.5$), and $\mathrm{pop}_j$ is cumulative popularity. Upon exposure, user reaction branches across a dual-valence affective spectrum. A positive reaction (like or affirmation, $r_{ij} = +1$) is triggered when attitudinal distance satisfies $d(\mathbf{x}_i, \mathbf{y}_j) \le \epsilon_{\mathrm{like}, i} \equiv \epsilon_i (1 + 0.30 e_j)$. The user assimilates toward $\mathbf{y}_j$,
\begin{equation}
    \Delta \mathbf{x}_i^{\mathrm{like}}(t) = \kappa_{\mathrm{feed}}\,\alpha_i\,(1 + 0.45 e_j)\,\exp\!\big(0.30\,b_i \max(0, \cos\theta_{ij})\big) \big(\mathbf{y}_j - \mathbf{x}_i\big),
    \label{eq:like_assimilation}
\end{equation}
incrementing the post like counter $\mathrm{Likes}_j \leftarrow \mathrm{Likes}_j + 1$. Conversely, a negative reaction (dislike, outrage, or quote-dunking, $r_{ij} = -1$) is triggered when content is strongly discordant, satisfying $d(\mathbf{x}_i, \mathbf{y}_j) \ge \tau_{\mathrm{dislike}, i} \equiv \max(\epsilon_{\mathrm{like}, i} + 0.08, 1.45\,\epsilon_i (1 - 0.18 e_j))$. Out-group indignation probability scales with discordance and confirmation bias $b_i$,
\begin{equation}
    P_{\mathrm{dislike}} = \min\!\Big(0.85,\, \kappa_{\mathrm{dislike}}(0.28 + 0.52 e_j)(0.45 + 0.35\min(2, b_i))\big(0.40 + 0.60(d - \tau_{\mathrm{dislike}})\big)\Big).
    \label{eq:dislike_prob}
\end{equation}
When triggered, the user registers out-group indignation ($\mathrm{Dislikes}_j \leftarrow \mathrm{Dislikes}_j + 1$) and undergoes repulsive backfire differentiation away from the post \citep{bail2018exposure},
\begin{equation}
    \Delta \mathbf{x}_i^{\mathrm{dislike}}(t) = \kappa_{\mathrm{feed}}\,\kappa_{\mathrm{backfire}}\,\alpha_i\,(0.35 + 0.65 e_j) \min(1.8, 0.5 + 0.35 b_i) \big(\mathbf{x}_i - \mathbf{y}_j\big).
    \label{eq:dislike_repulsion}
\end{equation}

\paragraph{Algorithmic Message Relevance and Controversy Amplification.}
Crucially, recommendation engines do not penalize negative reactions; rather, aggregate message relevance $R_m$ compounds both positive and negative engagements,
\begin{equation}
    R_m = \mathrm{Likes}_m + \omega \cdot \mathrm{Dislikes}_m, \qquad \omega \approx 1.0.
    \label{eq:message_relevance_methods}
\end{equation}
Because both in-group admiration and out-group indignation increase $R_m$, content situated along the diagonal of viral controversy achieves maximal algorithmic amplification (Fig.~\ref{fig:figure3}c) \citep{crockett2017moral}.

\paragraph{Multi-Stage Meme Diffusion and Dark-Social Cascades.}
Meme diffusion progresses across three distinct sequential phases \citep{leskovec2009meme}. In the algorithmic discovery stage, content is surfaced via recommendation feeds (Instagram and TikTok), driving rapid initial engagement among high-plasticity youth cohorts within $\tau_1 \approx 10$ steps. During the public debate and remixing stage, the meme permeates microblogging and commentary channels (YouTube and X), where quote-dunking and viral debate sustain peak controversy ($\tau_2 \approx 15\text{--}25$ steps). Finally, in the closed messaging cascades stage, narratives spill into private groups (WhatsApp and Telegram), where end-to-end encryption shields content from centralized moderation and high interpersonal trust ($\tau_{\mathrm{WA}} = 0.85$) sustains viral peer-to-peer forwarding over long horizons ($t > 30$ steps) \citep{resende2019whatsapp, garimella2020images}.

\paragraph{Hybrid Cross-Media Coupling.}
The traditional media and digital ecosystems interact through asymmetric agenda-setting loops \citep{chadwick2017hybrid, harder2017intermedia},
\begin{equation}
    \tau_{\mathrm{turnaround}} \gg \tau_{\mathrm{social}},
    \label{eq:timescale_ratio}
\end{equation}
where legacy newsrooms provide authoritative investigative releases with high initial credibility, and social platforms respond rapidly with remixing and partisan debate. In turn, aggregate digital volume $\mathcal{V}_{\mathrm{soc}}(t) = \sum_m R_m(t)$ exerts intermedia agenda-setting pressure back upon traditional outlets, compelling broadcast television and the press to maintain extended follow-up reporting.

\begin{table}[htbp]
\centering
\caption{Summary of traditional mass media and algorithmic social media platform parameters.}
\label{tab:media_parameters}
\footnotesize
\setlength{\tabcolsep}{4.5pt}
\begin{tabular}{lp{5.1cm}p{3.6cm}p{3.8cm}}
\toprule
\textbf{Symbol} & \textbf{Description} & \textbf{Value / Range} & \textbf{Reference} \\
\midrule
\multicolumn{4}{l}{\textit{Traditional Mass Media Ecosystem}} \\
$M, \mathbf{m}_g, \tau_g$ & Media groups, ideological centers, tendentiousness & $M \in \{0, 1, 2, 3, 4, 12\}, \tau_g \in [0, 1]$ & \citet{mullainathan2005market, gentzkow2006media} \\
$v \in \{0, 1, 2, 3, 4\}$ & Modalities (TV, Radio, Press, Outdoor, Rally) & 5 channels per conglomerate & \citet{kantar2023inside, reuters2023digital} \\
$\gamma_v, \delta_v$ & Modality credibility and attention decay rate & $\gamma \in [0.45, 0.82], \delta \in [0.08, 0.70]$ & \citet{wu2007novelty, downs1972up} \\
$\beta_m, \kappa_0$ & Media selective exposure beta and broadcast strength & $\beta_m \in [0.15, 2.16], \kappa_0 = 0.05$ & \citet{dandekar2013biased, dellavigna2007fox} \\
\midrule
\multicolumn{4}{l}{\textit{Algorithmic Social Media Platforms \& Influencers}} \\
$p \in \{0, 1, 2, 3\}$ & Platforms (Instagram/TikTok, YouTube/X, FB, WhatsApp) & 4 platform modalities & \citet{reuters2023digital, cinelli2021echo} \\
$\alpha, \mathrm{Gini}$ & Influencer follower power-law exponent and Gini & $\alpha \approx 1.15, \mathrm{Gini} \approx 0.73$ & \citet{barabasi1999emergence, bianconi2001competition} \\
$\beta, \mu_{\mathrm{emo}}$ & Algorithmic curation homophily and emotional boost & $\beta = 2.0, \mu_{\mathrm{emo}} = 1.5$ & \citet{brady2017emotion, vosoughi2018spread} \\
$\omega, R_m$ & Dual-valence engagement weight and relevance & $\omega = 1.0, R_m = \mathrm{Likes} + \omega\,\mathrm{Dislikes}$ & \citet{crockett2017moral, bail2018exposure} \\
$K_{\mathrm{WA}}, \tau_{\mathrm{WA}}$ & WhatsApp group capacity limit and peer trust weight & $K \le 256, \tau_{\mathrm{WA}} = 0.85$ & \citet{resende2019whatsapp, garimella2020images} \\
$\tau_{\mathrm{inf}}$ & Influencer creator tendentiousness & $[0.10, 0.90]$ & \citet{sasahara2021social, tucker2018social} \\
\bottomrule
\end{tabular}
\end{table}

% ------------------------------------------------------------------------------
\subsection{Electoral Module}
\label{subsec:electoral_module}

To evaluate how the micro-level attitudinal trajectories and conviction dynamics established across our three communicative substrates, namely the multiplex interpersonal network $\mathcal{M}$ (Section~\ref{subsec:multiplex_network}, Fig.~\ref{fig:figure1}), traditional mass media conglomerates (Section~\ref{subsec:methods_traditional_media}, Fig.~\ref{fig:figure2}), and algorithmic social media platforms (Section~\ref{subsec:methods_social_media}, Fig.~\ref{fig:figure3}), translate into macro-electoral outcomes, we formulate an electoral decision module grounded in dual affective evaluations, candidate affinity or likelihood ($A_{ik} \in [0, 1]$) and candidate aversion or rejection ($R_{ik} \in [0, 1]$) \citep{soroka2014negativity, lau2007effects}. In a multi-candidate electoral field (such as the First Leg contest among Lula $\mathbf{c}_L$, Fl{\'a}vio Bolsonaro $\mathbf{c}_F$, and Third-Way centrist contenders $\mathbf{c}_O$), each political contender is defined by a fixed coordinate centroid $\mathbf{c}_k \in \Omega = [-1, 1]^D$ within the identical four-dimensional metric concept space defined in Section~\ref{subsec:concept_space} (Fig.~\ref{fig:figure1}b). At each discrete simulation step $t$, a citizen's dynamic ideological distance to candidate $k$ is $d_{ik}(t) = \frac{1}{\sqrt{D}}\|\mathbf{x}_i(t) - \mathbf{c}_k\|_2$, which directly reflects the continuous integration of interpersonal persuasion (Eq.~\eqref{eq:fj_interpersonal_update}), broadcast exposure (Eq.~\eqref{eq:traditional_opinion_shift}), and digital algorithmic feeds (Eqs.~\eqref{eq:like_assimilation}--\eqref{eq:dislike_repulsion}).

\paragraph{Dynamic Evolution of Candidate Likelihood and Aversion.}
Candidate likelihood $A_{ik}(t)$ and aversion $R_{ik}(t)$ evolve as separate cognitive state variables with asymmetric update thresholds and memory decay rates,
\begin{align}
    A_{ik}(t+1) &= (1 - \delta_A)\,A_{ik}(t) + \lambda_A \max\!\left(0, \frac{\theta_A - d_{ik}(t)}{\theta_A}\right)^2, \label{eq:adm_dynamics} \\
    R_{ik}(t+1) &= (1 - \delta_R)\,R_{ik}(t) + \lambda_R \min\!\left(1, \max\!\left(0, \frac{d_{ik}(t) - \theta_R}{1.35 - \theta_R}\right)^{1.35}\right), \label{eq:rej_dynamics}
\end{align}
where $\theta_A = 0.62$ defines the threshold for positive ideological affinity, $\theta_R = 0.62$ is the threshold of ideological dissonance, $\lambda_A = 0.12$ and $\lambda_R = 0.14$ denote acquisition rates, and $\delta_A = 0.06$ and $\delta_R = 0.04$ are passive memory decay rates. Crucially, because $\delta_R < \delta_A$ and $\lambda_R > \lambda_A$, negative evaluations form more rapidly and decay more slowly than positive evaluations, capturing the well-established psychological negativity bias in political decision-making \citep{soroka2014negativity, lau2007effects}.

\paragraph{Asymmetric Negative Campaigning and Rejection Shocks.}
In modern digital political warfare, campaigns strategically deploy targeted rejection news focusing on scandal framing, moral outrage, and character attacks \citep{janones2023janonismo, cesarino2020posverdade}. These asymmetric rejection impulses are disseminated through the hybrid digital architecture analyzed in Section~\ref{subsec:methods_social_media}, where tendentious influencers (Fig.~\ref{fig:figure3}b) exploit outrage-weighted algorithmic discovery ($R_m$, Eq.~\eqref{eq:algorithmic_relevance}, Fig.~\ref{fig:figure3}c) and closed WhatsApp forwarding cascades (Fig.~\ref{fig:figure3}d), complemented by traditional newsroom reverberations (Fig.~\ref{fig:figure2}c, Fig.~\ref{fig:figure3}e). When citizen $i$ is exposed to a rejection news impulse targeting candidate $k$, their aversion is augmented directly,
\begin{equation}
    R_{ik}(t) \leftarrow \min\!\big(1.0,\, R_{ik}(t) + \lambda_{\mathrm{rej}} \cdot \mathbb{I}_{\mathrm{exposed}}(i)\big),
    \label{eq:rejection_shock}
\end{equation}
where $\lambda_{\mathrm{rej}}$ scales with campaign intensity. In our diagnostic scenario (Fig.~\ref{fig:figure4}), the pre-election sprint ($t \ge 80$) introduces an intense viral blitz ($\lambda_{\mathrm{rej}} \approx 0.065 \times [1 + 0.06(t - 80)]$, reaching $70\%$ of the electorate through closed dark-social cascades). Because $\delta_R < \delta_A$ and $\lambda_R > \lambda_A$ in Eqs.~\eqref{eq:adm_dynamics}--\eqref{eq:rej_dynamics}, negative rejection impulses induce enduring attitudinal scars that decay much more slowly than positive policy achievements.

\paragraph{Candidate Utility and Three-Way Voting Calculus.}
On election day, citizen $i$ evaluates the overall utility of each candidate $k$,
\begin{equation}
    U_{ik}(t) = \beta_A A_{ik}(t) - \beta_R R_{ik}(t) - \beta_d d_{ik}(t) + \xi_{ik}(t),
    \label{eq:candidate_utility}
\end{equation}
with weights $\beta_A = 1.15$, $\beta_R = 1.45$ reflecting the primacy of aversion over affinity, $\beta_d = 0.45$, and idiosyncratic evaluation noise $\xi_{ik} \sim \mathcal{N}(0, \sigma_{\mathrm{eval}}^2)$ ($\sigma_{\mathrm{eval}} = 0.03$). Citizens resolve their ballot through a three-way decision rule partitioning the electorate into Abstention ($C_i = -1$), Blank or Null Protest ($C_i = -2$), and Valid Vote ($C_i = k \in \{0, 1\}$). Under Abstention ($C_i = -1$), citizens refrain from voting due to political alienation when no candidate generates sufficient enthusiasm ($\max_k A_{ik} < \tau_{\mathrm{turnout}} = 0.14$) and no intense singular rejection compels a defensive tactical vote, or when mutual rejection is pervasive ($\min_k R_{ik} > 0.60$). Under Blank or Null Protest ($C_i = -2$), citizens turn out to cast a blank or spoiled protest ballot when the candidate utility margin is negligible ($|U_{i0} - U_{i1}| < \theta_{\mathrm{margin}} = 0.10$) and neither candidate triggers extreme aversion ($\max_k R_{ik} < \tau_{\mathrm{strong\_rej}} = 0.52$). Otherwise, the citizen casts a Valid Choice ($C_i = k$) for their highest-utility candidate, $C_i = \arg\max_k U_{ik}$.

% ------------------------------------------------------------------------------
\subsection{Algorithmic Counter-Offensive Framework: The Janones Tactical Doctrine}
\label{subsec:janones_tactical_doctrine}

To model how an asymmetric digital and institutional disadvantage can be actively neutralized, we formulate the tactical doctrine of \emph{Janonismo Cultural} as an explicit set of behavioral rules, state transitions, and algorithmic mechanics \citep{janones2023janonismo}. Rather than operating as disconnected heuristics, each of the six canonical Janones counter-tactics represents an intentional operationalization or structural inversion of the baseline communication mechanisms modeled in Sections~\ref{subsec:multiplex_network}--\ref{subsec:methods_social_media} (Figs.~\ref{fig:figure1}--\ref{fig:figure3}).

\paragraph{T1: Algorithmic Asphyxiation and Autonomous Agenda Setting.}
In standard online discourse, partisan counter-speech inadvertently amplifies opponent content because recommendation algorithms optimize for gross engagement ($R_m = \mathrm{Likes}_m + \omega \cdot \mathrm{Dislikes}_m$, Eq.~\eqref{eq:algorithmic_relevance}) \citep{crockett2017moral, bail2018exposure} (Fig.~\ref{fig:figure3}c). Under the Janones doctrine, allied nodes enforce a strict behavioral embargo on rival posts,
\begin{equation}
    P_{\mathrm{reply}}(i \to j) = 0, \quad P_{\mathrm{quote}}(i \to j) = 0, \quad P_{\mathrm{repost}}(i \to j) = 0, \qquad \forall i \in \mathrm{Allied}, \, j \in \mathrm{Rival}.
    \label{eq:asphyxiation_rule}
\end{equation}
By setting cross-engagement probabilities to zero, rival posts are deprived of the $\omega \cdot \mathrm{Dislikes}_m$ friction signals required by recommendation engines to trigger cross-cluster feed injection, establishing an attention starvation zone (Fig.~\ref{fig:figure5}b). Concurrently, allied nodes refuse reactive refutation, exploiting the separation of timescales where digital reactions outpace legacy gatekeeping ($\tau_{\mathrm{social}} \ll \tau_{\mathrm{turnaround}}$, Eq.~\eqref{eq:timescale_ratio}). When rival content exceeds an attention threshold $\theta_{\mathrm{att}}$, allied nodes inject an orthogonal, high-arousal counter-narrative $\mathbf{y}_{\mathrm{inject}} \in \Omega$ satisfying $\cos\theta(\mathbf{y}_{\mathrm{inject}}, \mathbf{y}_{\mathrm{rival}}) \approx 0$. Because the public attention budget is strictly finite ($\sum_k A_k(t) = 1.0$), saturating open feeds with an autonomous counter-agenda directly cannibalizes the viral headroom of rival topics \citep{wu2007novelty}.

\paragraph{T2: Strict Monothematism and Quadratic Algorithmic Virality.}
Rather than dispersing communication bandwidth across disparate policy topics, allied networks synchronize around a single topic $T_{\mathrm{day}}$ for each discrete 24-hour cycle,
\begin{equation}
    \sigma^2(\{T_i(t)\}_{i \in \mathrm{Allied}}) \to 0, \qquad \forall t \in [t_0, t_0 + 24].
    \label{eq:monothematism}
\end{equation}
The mathematical foundation of monothematism rests on algorithmic convexity inherent to platform recommendation probability functions (Eq.~\eqref{eq:recom_prob}) (Fig.~\ref{fig:figure5}a). In recommendation and trending algorithms, viral discovery scales superlinearly with message density, $V(\mathbf{x}) \propto \sum_{k=1}^K x_k^2$. For a fixed total communicative bandwidth $S = \sum_{k=1}^K x_k = 4$, distributing attention evenly across four issues yields $\sum x_k^2 = 1^2 + 1^2 + 1^2 + 1^2 = 4$, splitting across two topics yields $2^2 + 2^2 = 8$, while concentrating bandwidth onto a single monothematic agenda yields $4^2 = 16$. Because $4 < 8 < 16$, monothematic concentration doubles the effective virality of split messaging and quadruples that of dispersed communication, ensuring that allied content reliably breaches trending thresholds.

\paragraph{T3: Material Loss-Aversion Framing.}
Abstract institutional and democratic norms resonate weakly among lower-SES and undecided strata \citep{zaller1992nature}. The Janones doctrine substitutes abstract civic discourse with immediate, visceral, and tangible material threats, such as minimum wage de-indexation, working-hour extensions, and welfare cuts. In the opinion and utility updates of undecided citizens, economic threat inputs framed around prospective losses ($\Delta U < 0$) are weighted by a Kahneman--Tversky loss multiplier $\lambda_{\mathrm{loss}} \in [2.0, 2.5]$ \citep{kahneman1979prospect},
\begin{equation}
    \Delta \mathbf{x}_i^{\mathrm{loss}}(t) = \lambda_{\mathrm{loss}}\,\kappa_{\mathrm{feed}}\,\alpha_i\,\mathrm{SES}_i^{-0.35} \big(\mathbf{y}_{\mathrm{loss}} - \mathbf{x}_i(t)\big),
    \label{eq:loss_aversion_update}
\end{equation}
which modifies the interpersonal attitude trajectory (Eq.~\eqref{eq:fj_interpersonal_update}) and overcomes the slow inertia of conviction anchoring $\mathbf{s}_i$ (Eq.~\eqref{eq:conviction_update}, Fig.~\ref{fig:figure1}d), mobilizing economically vulnerable citizens who are otherwise prone to electoral abstention.

\paragraph{T4: Molecular WhatsApp Synchronization and Phased Dispatch.}
Content assets are formatted as raw, informal, unsponsored video clips and voice notes that strip institutional branding. Deliberative debate is suspended during the active campaign while a centralized command hub pushes scripted assets to high-degree subnetworks in three coordinated daily phases across the multi-platform architecture (Fig.~\ref{fig:figure3}a), spanning morning closed-network morale alignment, afternoon open-feed hashtag surges, and evening molecular relays into private family and local community networks. Allied volunteers relay these units from broadcast channels directly into intimate, natural family and neighborhood groups ($\mathcal{L}_3^{\mathrm{soc}}$). Because interpersonal trust in closed groups is exceptionally high ($\tau_{\mathrm{WA}} = 0.85$, Fig.~\ref{fig:figure3}d), information transmission exhibits a conversion efficiency up to $2.8\times$ higher than open algorithmic feeds \citep{resende2019whatsapp}, effectively mapping onto the high-weight edges of the physical multiplex layers $\mathcal{L}_1$ (household) and $\mathcal{L}_4$ (friends) defined in Section~\ref{subsec:multiplex_network} (Fig.~\ref{fig:figure1}a).

\paragraph{T5: Mass In-Person Capital Rallies and Physical Turnout Locking.}
To translate digital enthusiasm into physical electoral action, the campaign organizes open-air rallies across major metropolitan capitals, including S{\~a}o Paulo, Rio de Janeiro, Belo Horizonte, Salvador, Recife, and Porto Alegre. Rallies operationalize the traditional media public demonstration modality ($v=4$, Eq.~\eqref{eq:reach_rally}, Section~\ref{subsec:methods_traditional_media}, Fig.~\ref{fig:figure2}a) and directly activate the multiplex physical layers ($\mathcal{L}_1, \mathcal{L}_2, \mathcal{L}_4, \mathcal{L}_5$, Fig.~\ref{fig:figure1}a). Assembling hundreds of thousands of citizens in public thoroughfares generates three coupled systemic effects, namely physical turnout locking where collective effervescence lifts candidate admiration above the turnout threshold ($A_{i,L} > \tau_{\mathrm{turnout}}$), interpersonal contagion multipliers as attendees return as high-conviction opinion leaders radiating enthusiasm through high-trust household ($\mathcal{L}_1$) and friendship ($\mathcal{L}_4$) ties, and hybrid media feedback where massive aerial crowd footage provides authentic visual assets that dominate digital feeds (Fig.~\ref{fig:figure3}e).

\paragraph{T6: The 24-Hour Live Broadcast and Signal Saturation.}
During the final pre-election sprint ($t \ge 80$), a continuous 24-hour live broadcast event saturates the public information ecosystem with high-interest discussions. By maintaining an unbroken, high-salience communicative ceiling across the hybrid media ecosystem (Fig.~\ref{fig:figure3}e), the live broadcast absorbs $62\text{--}65\%$ of public attention, preventing the opponent blitz from securing uncontested viral dominance and damping rival reach by $65\%$ (Fig.~\ref{fig:figure5}f). Table~\ref{tab:janones_tactics} summarizes the complete operationalization of the Janones doctrine into agent-based graph mechanics.

\begin{table}[htbp]
\centering
\caption{Tactical translation table: operationalizing the Janones doctrine into agent-based graph mechanics across the foundational model layers.}
\label{tab:janones_tactics}
\footnotesize
\setlength{\tabcolsep}{4pt}
\begin{tabular}{lp{4.2cm}p{3.8cm}p{2.8cm}}
\toprule
\textbf{Janones Tactic} & \textbf{Operational Description} & \textbf{Agent / Graph Mechanism} & \textbf{Parameter / Subsystem} \\
\midrule
T1 Asphyxiation & Suppression of outbound interaction toward rival bait & Interactivity probability $P_{\mathrm{reply}}, P_{\mathrm{quote}} \to 0$ & $R_m$ signal denial (Fig.~\ref{fig:figure3}c) \\
T2 Monothematism & One coordinated agenda per 24h cycle & Variance $\sigma^2(T_{\mathrm{day}}) \to 0$ across allied nodes & Quadratic virality $\sum x_k^2$ (Fig.~\ref{fig:figure5}a) \\
T3 Material Loss Frame & Concrete economic risks (welfare cuts, wage freeze) & Kahneman--Tversky loss weighting & $\lambda_{\mathrm{loss}} \in [2.0, 2.5]$ (Fig.~\ref{fig:figure1}d) \\
T4 WhatsApp Sync & Scripted asset dispatch and molecular relay & Phased dispatch into private cliques & Trust weight $\times 2.8$ (Fig.~\ref{fig:figure3}d, $\mathcal{L}_1/\mathcal{L}_4$) \\
T5 Mass Rallies & Vast public assemblies in metropolitan capitals & Modality $v=4$ and multiplex layers $\mathcal{L}_1/\mathcal{L}_4$ & Turnout locking $A_{i,L} > \tau_{\mathrm{turnout}}$ \\
T6 24h Live Saturation & Nonstop continuous broadcast during pre-election sprint & Attention saturation ceiling & Capping rival blitz by $65\%$ (Fig.~\ref{fig:figure3}e) \\
\bottomrule
\end{tabular}
\end{table}

\subsection{Scaling Across City Population Sizes}
\label{subsec:city_scaling}

To investigate how urban population scale affects emergent attitude distributions and cross-issue alignment in a pure interpersonal network, we simulate four city scales spanning four orders of magnitude \citep{bettencourt2007growth, schlapfer2014scaling}, comprising Small City (\texttt{small\_10k}, $N = 10^4$), Medium City (\texttt{medium\_100k}, $N = 10^5$, reference city scale for deep multi-timescale experiments), Large City (\texttt{large\_1m}, $N = 10^6$), and Metropolis (\texttt{metropolis\_10m}, $N = 10^7$). Across all four scales, the intensive demographic distributions (Eqs.~\eqref{eq:age_norm}--\eqref{eq:assimilation_bias}) and the mean degree of each interpersonal network layer ($\langle k^{[1]} \rangle \approx 1.8$, $\langle k^{[2]} \rangle = 10$, $\langle k^{[3]} \rangle = 8$, $\langle k^{[4]} \rangle = 6$, $\langle k^{[5]} \rangle = 4$) are held strictly constant. Consequently, while a citizen's local ego-network size remains invariant with respect to $N$, the global network diameter, the number of distinct household cliques ($\approx N / 2.79$), and the diversity of cross-cutting weak ties ($\mathcal{E}^{[5]}$) scale linearly with $N$. In smaller cities ($N = 10^4$), finite-size topological clustering induces stronger shared local fluctuations across issue dimensions, whereas in larger cities ($N \ge 10^5$), independent local neighborhoods self-average across the population \citep{castellano2009statistical}.

% ------------------------------------------------------------------------------
\subsection{Exogenous Fact Interventions}
\label{subsec:interventions}

To test the dynamical separation of timescales between fast current opinions $\mathbf{x}_i(t)$ and slow convictions $\mathbf{s}_i(t)$, we design two controlled exogenous fact interventions in the medium city ($N = 10^5$) without altering interpersonal network rules:

\paragraph{1. Temporary Single High-Salience Event Intervention.}
To measure how a single important real-world event perturbs public opinion while leaving underlying convictions largely intact, the population evolves under a neutral fact environment ($\mathbf{b}_{\mathrm{fact}}(t) = \mathbf{0}$, $I_{\mathrm{fact}}(t) = 0$) during a baseline stabilization period $0 \le t < t_{\mathrm{start}}$ ($t_{\mathrm{start}} = 450$). During a bounded event window $t \in [t_{\mathrm{start}}, t_{\mathrm{end}}] = [450, 570]$ of duration $\Delta t_{\mathrm{ev}} = 120$ steps, a high-salience exogenous fact vector is injected into the concept space:
\begin{equation}
    \mathbf{b}_{\mathrm{fact}}(t) =
    \begin{cases}
        (+0.85, +0.60, -0.50, +0.40)^\top, & \text{for } 450 \le t \le 570 \quad (I_{\mathrm{fact}}(t) = 1), \\[4pt]
        (0, 0, 0, 0)^\top, & \text{for } t < 450 \text{ or } t > 570 \quad (I_{\mathrm{fact}}(t) = 0),
    \end{cases}
    \label{eq:single_event_intervention}
\end{equation}
with event pull $\gamma_0 = 0.075$. At $t > 570$, the temporary event ceases and the environment returns to neutral baseline facts ($\mathbf{b}_{\mathrm{fact}}(t) = \mathbf{0}$), allowing direct measurement of (i) the peak excursion ratio $|\Delta \bar{x}_0(t_{\mathrm{end}})| / |\Delta \bar{s}_0(t_{\mathrm{end}})|$ and (ii) the post-event relaxation of $\mathbf{x}_i(t)$ back toward the conviction anchor $\mathbf{s}_i(t)$.

\paragraph{2. Multi-Phase Fact-Bias Regime Shift Intervention.}
To evaluate long-horizon resilience and plastic drift under sustained macro-environmental shocks over $T = 1500$ steps, we implement a three-phase regime-shift protocol (\texttt{regime\_shift}):
\begin{equation}
    b_{\mathrm{fact}, d}(t) = (-1)^d A_{\mathrm{fact}}(t), \qquad
    A_{\mathrm{fact}}(t) =
    \begin{cases}
        0.00, & 0 \le t \le 400 \quad \text{(Phase I: Neutral Crystallization)}, \\[4pt]
        +0.78, & 401 \le t \le 950 \quad \text{(Phase II: Positive Fact-Bias Shift)}, \\[4pt]
        -0.82, & 951 \le t \le 1500 \quad \text{(Phase III: Negative Fact-Bias Reversal)}.
    \end{cases}
    \label{eq:regime_shift_intervention}
\end{equation}

% ------------------------------------------------------------------------------
\subsection{Evaluation Metrics and Statistical Analysis}
\label{subsec:metrics_analysis}

At regular snapshot intervals across every simulation trajectory, we compute a suite of aggregate and cohort-stratified metrics to characterize central tendency, dispersion, multi-timescale volatility, clustering, and multidimensional issue alignment.

\paragraph{1. Mean Normalized Distance to the Center and Attitude Extremity.}
To track convergence toward a centrist distribution in $\Omega = [-1, 1]^D$, we compute the population-average normalized Euclidean distance to the origin $\mathbf{0}$ separately for convictions and current opinions \citep{howe2017attitude}:
\begin{equation}
    \bar{E}_{\mathbf{s}}(t) = \frac{1}{N}\sum_{i=1}^N \frac{\|\mathbf{s}_i(t)\|_2}{\sqrt{D}}, \qquad
    \bar{E}_{\mathbf{x}}(t) = \frac{1}{N}\sum_{i=1}^N \frac{\|\mathbf{x}_i(t)\|_2}{\sqrt{D}} \in [0, 1].
    \label{eq:mean_distance_to_center}
\end{equation}

\paragraph{2. Opinion--Conviction Excursion Gap and Step-to-Step Volatilities.}
The instantaneous decoupling between citizens' current opinions and their underlying convictions is measured by the mean opinion--conviction excursion gap $G(t)$:
\begin{equation}
    G(t) = \frac{1}{N}\sum_{i=1}^N G_i(t), \qquad G_i(t) = \frac{1}{\sqrt{D}} \|\mathbf{x}_i(t) - \mathbf{s}_i(t)\|_2.
    \label{eq:op_conv_gap}
\end{equation}
To compare the rate of temporal fluctuation on the fast and slow timescales, we compute the step-to-step opinion volatility $V_{\mathbf{x}}(t)$ and conviction volatility $V_{\mathbf{s}}(t)$:
\begin{equation}
    V_{\mathbf{x}}(t) = \frac{1}{N}\sum_{i=1}^N \frac{\|\mathbf{x}_i(t) - \mathbf{x}_i(t-1)\|_2}{\sqrt{D}}, \qquad
    V_{\mathbf{s}}(t) = \frac{1}{N}\sum_{i=1}^N \frac{\|\mathbf{s}_i(t) - \mathbf{s}_i(t-1)\|_2}{\sqrt{D}}.
    \label{eq:step_volatilities}
\end{equation}
Both $G(t)$ and $V_{\mathbf{x}}(t), V_{\mathbf{s}}(t)$ are disaggregated across age cohorts (Young $<30$, Mid-Age $30\text{--}54$, Senior $\ge 55$), socioeconomic strata (Low, Mid, High SES), and conviction extremity tiers (Centrist, Moderate, Extremist).

\paragraph{3. Normalized Multidimensional Dispersion Ratio.}
Following \citet{bramson2016disambiguation}, population spread across the $D$ dimensions is benchmarked against the variance of a continuous uniform distribution on $[-1, 1]$, $\mathrm{Var}(\mathcal{U}[-1, 1]) = 1/3$:
\begin{equation}
    D_{\mathrm{ratio}}^{\mathbf{x}}(t) = \frac{\frac{1}{D}\sum_{d=0}^{D-1} \mathrm{Var}_i\!\big(x_{i,d}(t)\big)}{1/3} = \frac{3}{D}\sum_{d=0}^{D-1} \mathrm{Var}_i\!\big(x_{i,d}(t)\big),
    \label{eq:dispersion_ratio}
\end{equation}
and analogously $D_{\mathrm{ratio}}^{\mathbf{s}}(t) = \frac{3}{D}\sum_{d=0}^{D-1} \mathrm{Var}_i(s_{i,d}(t))$ for convictions. A value of $D_{\mathrm{ratio}} = 1$ corresponds to uniform random dispersion, $D_{\mathrm{ratio}} \to 0$ indicates full consensus, and $0.25 \le D_{\mathrm{ratio}} \le 0.40$ characterizes a cohesive centrist distribution.

\paragraph{4. Cross-Issue Alignment and Effective Ideological Dimensionality.}
In multidimensional attitude spaces, ideological sorting manifests not only within individual issues but through \textbf{issue alignment} (or constraint), defined as the degree to which a citizen's stance on one policy dimension predicts their stances on logically independent dimensions \citep{converse1964nature, baldassarri2008partisans}. Let $\boldsymbol{\Sigma}(t) \in \mathbb{R}^{D \times D}$ denote the empirical covariance matrix of current opinions $\{\mathbf{x}_i(t)\}_{i=1}^N$ with entries $\Sigma_{dd'}(t) = \mathrm{Cov}(x_{\cdot, d}(t), x_{\cdot, d'}(t))$, and let $\mathbf{R}(t) = [\rho_{dd'}(t)]$ denote the corresponding Pearson correlation matrix,
\begin{equation}
    \rho_{dd'}(t) = \frac{\Sigma_{dd'}(t)}{\sqrt{\Sigma_{dd}(t)\,\Sigma_{d'd'}(t)}} \in [-1, 1].
    \label{eq:pearson_issue_corr}
\end{equation}
Following \citet{baldassarri2008partisans}, we quantify cross-issue alignment using the mean and maximum absolute pairwise off-diagonal correlations across all $\binom{D}{2} = 6$ issue pairs:
\begin{equation}
    \bar{|\rho|}(t) = \frac{2}{D(D-1)} \sum_{0 \le d < d' \le D-1} |\rho_{dd'}(t)|, \qquad
    \rho_{\max}(t) = \max_{0 \le d < d' \le D-1} |\rho_{dd'}(t)|.
    \label{eq:issue_alignment_metrics}
\end{equation}
Complementarily, to measure whether the $D$-dimensional concept space remains multi-axis or collapses onto a single left--right ideological axis, we compute the \textbf{Effective Ideological Dimensionality} via Roy's eigenvalue participation ratio \citep{roy1953useful}:
\begin{equation}
    \mathrm{PR}(\boldsymbol{\Sigma}(t)) = \frac{\big(\mathrm{Tr}\,\boldsymbol{\Sigma}(t)\big)^2}{\mathrm{Tr}\!\big(\boldsymbol{\Sigma}(t)^2\big)} = \frac{\left(\sum_{d=1}^D \lambda_d(t)\right)^2}{\sum_{d=1}^D \lambda_d(t)^2} \in [1, D],
    \label{eq:participation_ratio}
\end{equation}
where $\lambda_1(t), \dots, \lambda_D(t) \ge 0$ are the eigenvalues of $\boldsymbol{\Sigma}(t)$. When the $D = 4$ dimensions vary independently with equal variance, $\mathrm{PR}(\boldsymbol{\Sigma}) = 4.0$; when all issues collapse onto a single ideological axis, $\mathrm{PR}(\boldsymbol{\Sigma}) \to 1.0$.

\paragraph{5. Polarization, Multimodality, and Network Segregation Indices.}
To test whether attitudes fragment into antagonistic poles or contract into a unimodal centrist distribution, we evaluate four additional structural metrics. The Multidimensional Esteban--Ray Polarization Index $P_{\mathrm{ER}}(t)$ \citep{esteban1994measurement} quantizes $\Omega = [-1, 1]^D$ into $B^D = 4^4 = 256$ cells with nonempty cell masses $\pi_c(t)$ and centroids $\mathbf{y}_c(t)$ ($c = 1, \dots, K$),
\begin{equation}
    P_{\mathrm{ER}}(t) = \kappa_{\mathrm{ER}} \sum_{c=1}^K \sum_{c'=1}^K \pi_c(t)^{1 + \alpha_{\mathrm{ER}}} \pi_{c'}(t) \|\mathbf{y}_c(t) - \mathbf{y}_{c'}(t)\|_2,
    \label{eq:esteban_ray}
\end{equation}
with polarization sensitivity parameter $\alpha_{\mathrm{ER}} = 1.3$ and scaling constant $\kappa_{\mathrm{ER}} = 1.0$. Sarle's Bimodality Coefficient $\mathrm{BC}(t)$ \citep{pfister2013good} is averaged across the $D$ dimensions using standardized skewness $\gamma_d$ and excess kurtosis $\kappa_d$,
\begin{equation}
    \mathrm{BC}(t) = \frac{1}{D}\sum_{d=0}^{D-1} \frac{\gamma_d(t)^2 + 1}{\kappa_d(t) + \dfrac{3(N-1)^2}{(N-2)(N-3)}},
    \label{eq:sarle_bc}
\end{equation}
where $\mathrm{BC} > 5/9 \approx 0.555$ indicates bimodality or multimodality, $\mathrm{BC} < 5/9$ indicates a unimodal distribution, and $\mathrm{BC} = 1/3 \approx 0.333$ for a Gaussian distribution. The Effective Number of Opinion Clusters \citep{laakso1979effective, jost2006entropy} is computed via the Inverse Participation Ratio $N_{\mathrm{eff}}^{\mathrm{IPR}}(t) = \big(\sum_{c=1}^K \pi_c(t)^2\big)^{-1}$ and exponential Shannon entropy $N_{\mathrm{eff}}^{H}(t) = \exp\!\big(-\sum_{c=1}^K \pi_c(t)\ln\pi_c(t)\big)$. Finally, Network Ideological Segregation $S_{\mathrm{net}}(t) \in [0, 1]$ is defined as one minus the ratio of mean weighted neighbor attitude distance to the expected distance of random pairs in the population,
\begin{equation}
    S_{\mathrm{net}}(t) = \mathrm{clip}_{[0, 1]}\!\left(1 - \frac{\frac{1}{N}\sum_{i=1}^N \sum_{j \in \mathcal{N}_i} W_{ij}\, d(\mathbf{x}_i(t), \mathbf{x}_j(t))}{\sqrt{\frac{2}{D}\sum_{d=0}^{D-1} \mathrm{Var}_i(x_{i,d}(t))}}\right).
    \label{eq:network_segregation}
\end{equation}

\paragraph{6. Marginal Attitude Density Estimation and Ensemble Statistical Aggregation.}
To inspect the full shape of the stationary attitude distributions at $t = T$, one-dimensional marginal probability densities $p(x_d)$ and $p(s_d)$ along each dimension $d \in \{0, 1, 2, 3\}$ are estimated via Gaussian kernel density estimation (KDE) with Scott's bandwidth rule $h = N^{-1/5}\hat{\sigma}_d$ \citep{scott2015multivariate}. Across all $S = 100$ Monte Carlo realizations for each scenario and city scale, time-series and terminal metrics are summarized by the ensemble mean $\bar{\mu}(t) = \frac{1}{S}\sum_{s=1}^S \mu^{(s)}(t)$, sample standard deviation $\hat{\sigma}_{\mu}(t)$, and $95\%$ confidence intervals $\bar{\mu}(t) \pm 1.96\,\hat{\sigma}_{\mu}(t)/\sqrt{S}$.




% ==============================================================================
% Results Section
% ==============================================================================

\section{Results}
\label{sec:results}


\subsection{The 2026 Brazilian Presidential Election}

Brazil operates as a presidential federative republic with concurrent general elections held every four years, allowing executive incumbents a single consecutive reelection alongside the renewal of the bicameral National Congress. In the 2026 contest, incumbent president Luiz In\'acio Lula da Silva of the Workers' Party (\emph{Partido dos Trabalhadores}, PT) sought reelection, challenged primarily by Senator Fl\'avio Bolsonaro of the right-wing Liberal Party (\emph{Partido Liberal}, PL). Fl\'avio is the eldest son of former president Jair Bolsonaro, who had narrowly lost the 2022 election to Lula and was subsequently incarcerated following judicial convictions related to an attempted coup d'\'etat. Third-party contenders exerted minimal electoral resonance, consisting largely of conservative figures attempting to capture the Bolsonarista electorate while Fl\'avio's ability to unify the right remained unproven.

The 2026 presidential race was defined by two sharply divergent campaign regimes across its first and second rounds. Throughout the opening 50-day campaign, voter intentions exhibited remarkable inertia (\textbf{Fig.~\ref{fig:election_2026_synthesis}b}): longitudinal tracking surveys by Datafolha depicted a steady, narrow lead for Lula (fluctuating between 41.8\% and 45.0\% of valid votes) over Fl\'avio (37.5\% to 41.0\%), with third-way candidates holding 14.0\% to 22.0\%. During this period, candidate rejection was mutually crystallized in a narrow band of 45\%--47\% (\textbf{Fig.~\ref{fig:election_2026_synthesis}c}). Neither front-runner participated in televised debates, and public discourse was largely overshadowed by controversies concerning prominent judiciary figures rather than substantive executive policy. 

This opening phase also reflected sharp structural asymmetries in institutional campaign endowments. Under Brazilian electoral law (Lei n\textordmasculine~9.504/1997), broadcast airtime on free television and radio (HGPE) is distributed proportionally to coalition representation in the Chamber of Deputies (\textbf{Fig.~\ref{fig:election_2026_synthesis}d}): Lula's broad seven-party coalition secured 5m~31s per daily block (44.1\% of total airtime; 433 30-second insertions) compared to Fl\'avio's 4m~20s (34.7\%; 339 insertions), yielding a $+27.3\%$ airtime advantage for the incumbent. Financially, while the PL controlled the largest national allocation from the public Special Campaign Financing Fund ($R\$$~881.7 million, 17.8\% of the $R\$$~4.96 billion global FEFC pool; \textbf{Fig.~\ref{fig:election_2026_synthesis}e}), the presidential campaigns operated under strict statutory spending limits ($R\$$~88.9 million in Round~1, expanding to $R\$$~133.4 million in Round~2; \textbf{Fig.~\ref{fig:election_2026_synthesis}e}). Crucially, the two tickets diverged in their revenue architecture (\textbf{Fig.~\ref{fig:election_2026_synthesis}f}): Lula deployed a centralized public apparatus funded 96.3\% by the FEFC ($R\$$~123.8M of $R\$$~128.6M declared), whereas Fl\'avio relied on a hybrid structure that offset lower party fund transfers ($R\$$~47.0M, 52.1\%) by raising an unprecedented $R\$$~43.2M in private individual contributions and crowdfunding, yielding a $9.0\times$ private capital advantage over Lula ($R\$$~4.8M).

The first-round ballot on October 4 produced a profound empirical shock (\textbf{Fig.~\ref{fig:election_2026_synthesis}a}). Fl\'avio Bolsonaro led with 47.03\% of the valid vote, finishing ahead of Lula (45.16\%) as third-party support collapsed to 7.81\%, while producing decisive legislative coattails for the PL in both congressional chambers. Furthermore, voter non-participation surged to a historical high of 24.70\% (20.90\% abstention and 4.80\% blank/null ballots; \textbf{Fig.~\ref{fig:election_2026_synthesis}a}). On election night, Bolsonarista social media mobilized aggressively, generating 3.90 million daily interactions compared to Lula's 0.17 million, representing a $22.9\times$ volume ratio (\textbf{Fig.~\ref{fig:election_2026_synthesis}h}).

This unexpected outcome prompted an immediate realignment. Initial post-election surveys projected a clear advantage for Fl\'avio (52.0\% vs.\ 48.0\% on October 8; \textbf{Fig.~\ref{fig:election_2026_synthesis}b}), supported by a post-victory honeymoon that lowered his rejection rate to 43.0\% while Lula's rose to 47.0\% (\textbf{Fig.~\ref{fig:election_2026_synthesis}c}), mirrored by a 9.3-point lead for Fl\'avio in the Quaest Digital Popularity Index (80.0 vs.\ 70.7; \textbf{Fig.~\ref{fig:election_2026_synthesis}g}). Under runoff rules, broadcast airtime was equalized to strict 5m~00s parity (25 spots/day each; \textbf{Fig.~\ref{fig:election_2026_synthesis}d}), shifting the decisive battleground toward the digital sphere.

In response, the second-round campaign adopted a fundamentally restructured tactical posture. Although Fl\'avio possessed a substantial $2.44\times$ multiplier in legislative surrogate bench strength (averaging 17.4 million cross-platform followers across his top-5 allies, led by Nikolas Ferreira with 52.0M, versus 7.1 million for PT allies; \textbf{Fig.~\ref{fig:election_2026_synthesis}i}), Lula counterbalanced this with a larger direct personal base (39.0M vs.\ 24.1M; \textbf{Fig.~\ref{fig:election_2026_synthesis}i}) and a massive closed WhatsApp infrastructure (\emph{Ex\'ercito do Lula}, 10.0 million participants vs.\ Fl\'avio's 2.0 million; \textbf{Fig.~\ref{fig:election_2026_synthesis}i}). Orchestrated around aggressive responsive messaging and material-loss framing regarding welfare protections, the Lulista digital counter-offensive dominated engagement during the October 16 televised debate (25.5M vs.\ 5.0M interactions; \textbf{Fig.~\ref{fig:election_2026_synthesis}h}) and sustained a $3.5\times$ daily interaction advantage throughout the runoff (8.52M vs.\ 2.41M/day; \textbf{Fig.~\ref{fig:election_2026_synthesis}h}).

This digital saturation steadily neutralized Fl\'avio's momentum. In public opinion, Fl\'avio's rejection rose from 43.0\% to 46.0\% by October 22 (\textbf{Fig.~\ref{fig:election_2026_synthesis}c}), while his polling advantage eroded from $+4.0\%$ to an absolute technical tie (50.0\% vs.\ 50.0\% on October 22; \textbf{Fig.~\ref{fig:election_2026_synthesis}b}). Simultaneously, the Quaest IPD inverted by $+16.0$ net points, swinging from Fl\'avio $+9.3$ to Lula $+6.7$ (83.3 vs.\ 76.6; \textbf{Fig.~\ref{fig:election_2026_synthesis}g}). Although the final election-eve Datafolha poll still placed Fl\'avio ahead within the margin of error (51.0\% vs.\ 49.0\%; \textbf{Fig.~\ref{fig:election_2026_synthesis}b}), election day witnessed an unprecedented democratic turnaround: universal free transit decrees and intensive mobilization compressed total non-voters by 2.91 percentage points (falling from 24.70\% to 21.79\%; \textbf{Fig.~\ref{fig:election_2026_synthesis}a}). This $+2.79$~pp net shift over eve polling culminated in a 51.79\% to 48.21\% victory for Lula, marking the first time in modern Brazilian history that a candidate trailing in the first round won the presidency in the runoff.


\begin{figure*}[t!]
\centering
  \includegraphics[width=\textwidth]{figures/figure0_jan_electiondata.pdf}%

\caption{\textbf{The 2026 Brazilian Presidential Election in Data.} 
\textbf{a}, Official TSE election results for Round~1 (Oct~4) and Round~2 (Oct~25), alongside turnout mobilization that reduced non-voters by $-2.91$~pp. 
\textbf{b}, Longitudinal tracking of valid vote intentions across six Round~1 and four Round~2 Datafolha polls, with gold markers indicating official ballot tallies. 
\textbf{c}, Candidate rejection dynamics (Datafolha: \emph{``would not vote for under any circumstances''}), highlighting the runoff reversal that increased Fl\'avio's rejection to 46\%. 
\textbf{d}, Statutory free broadcast airtime (HGPE) per daily block on television and radio, shifting from a $+27.3\%$ coalition advantage for Lula in Round~1 to statutory parity (5m~00s each) in Round~2. 
\textbf{e}, Macro public party financing (FEFC; total pool $R\$$~4.96B) and TSE presidential spending ceilings ($R\$$~88.9M in Round~1; $R\$$~133.4M cumulative in Round~2). 
\textbf{f}, Presidential ticket funding profiles, comparing Lula's state-financed budget (96.3\% FEFC) with Fl\'avio's hybrid apparatus, which has a $9.0\times$ private donation advantage. 
\textbf{g}, Quaest Digital Popularity Index (IPD; 0--100 scale, 175 metrics) across three milestone waves, showing a $+16.0$-point net turnaround swing for Lula. 
\textbf{h}, Daily social interaction volumes (millions/day; Bites data) during the Round~1 shock, televised debate, and the 13-day runoff window. 
\textbf{i}, Follower ecosystem and bench architecture, reporting candidate direct accounts across five platforms, official party handles, closed WhatsApp groups (Ex\'ercito do Lula 10.0M vs.\ Tropa do Fl\'avio 2.0M), and the arithmetic mean of the top-5 legislative surrogates per coalition.}
\label{fig:election_2026_synthesis}
\end{figure*}


\subsection{A Multi-Layer Digital Twin of the Electoral Battleground}

A defining feature of the 2026 presidential election is that the digital strategy deployed by Lula's runoff campaign had been publicly codified and advocated beforehand by Federal Deputy Andr\'e Janones. Premised on the diagnosis that Lula possessed broader underlying popular appeal but lacked an effective digital communication apparatus, Janones articulated his tactical doctrine for online political warfare in his book \emph{Janonismo Cultural}~\citep{janones2023janonismo}. Consequently, the 2026 Brazilian race offers a rare quasi-natural experiment in which two contrasting communication strategies were executed sequentially by the same candidate across the first and second rounds, backed by abundant longitudinal data that allows us to isolate and quantify the efficacy of specific tactical mechanisms. To this end, we develop a multi-layer network model of citizens whose political attitudes evolve under competing informational and social forces, providing a quantitative lens into how contemporary electoral politics is contested on the digital battleground. Drawing on modeling frameworks from computational epidemiology, we construct a calibrated digital twin of the Brazilian electoral environment. By fitting this model to the empirical trajectory of the 2026 race, we estimate the causal impact of the second-round strategic pivot. Our framework integrates three interacting influence layers: (i)~the urban interpersonal network, capturing everyday face-to-face interactions across households, workplaces, schools, public transit, religious congregations, and social activities; (ii)~the traditional media and ground-campaign layer, encompassing broadcast television, radio, print journalism, and political rallies; and (iii)~the digital media layer, which models algorithmic amplification and viral propagation across social networking platforms. The formal construction and parameterization of this multi-layer network are detailed below. Readers interested directly in the findings may skip to the next subsection (Subsection~\ref{subsec:results_electoral_janones}).



\subsubsection*{The Urban Interpersonal Attitude Network}

\begin{figure*}[ht!]
\centering
\includegraphics[width=\textwidth]{figures/figure1_jan_nonmedianetwork.pdf}
\caption{\textbf{Structure and dynamics of the urban interpersonal attitude network.} 
\textbf{a},~Empirical age pyramid and age-dependent cognitive plasticity $\alpha(a)$.
\textbf{b},~Urban socioeconomic status (SES) distribution.
\textbf{c},~Distributions of individual behavioral traits ($\alpha, \epsilon, u, b$).
\textbf{d},~Five-layer multiplex network topology $\mathcal{M} = \{\mathcal{L}_1, \dots, \mathcal{L}_5\}$.
\textbf{e},~Inter-cohort age contact matrix $M_{ab}$.
\textbf{f},~Age homophily $h_{\mathrm{age}}$ and ideological assortativity $r_{\mathrm{ideol}}$ across layers $\mathcal{L}_1$--$\mathcal{L}_5$.
\textbf{g},~Two-dimensional projection of initial ($\mathbf{s}_i(0)$) and converged ($\mathbf{s}_i(T)$) convictions alongside exogenous facts $\mathbf{f}_k$.
\textbf{h},~Temporal evolution of pairwise issue alignment $\langle |\rho_{kl}(t)| \rangle$ across city sizes ($N = 10\mathrm{k}$ to $10\mathrm{M}$).
\textbf{i},~Asymptotic alignment ($\max |\rho|$, $\langle |\rho| \rangle$) and effective dimensionality $\mathrm{PR}(\Sigma)$ versus urban scale.
\textbf{j},~Stationary marginal densities $p(s_d)$ across ideological dimensions $D_1$--$D_4$.
\textbf{l},~Two-timescale relaxation of mean opinion $\bar{\mathbf{x}}(t)$ and conviction $\bar{\mathbf{s}}(t)$ under a transient shock.
\textbf{m},~Age-stratified opinion--conviction tether stretching $\|\mathbf{x}_i - \mathbf{s}_i\|$.
\textbf{n},~Attitude space snapshot at peak shock ($t_{\mathrm{shock}}$) contrasting displaced opinions $\mathbf{x}_i$ and anchored convictions $\mathbf{s}_i$.}
\label{fig:figure1}
\end{figure*}

To establish the foundational layer of our electoral digital twin, we first characterize the baseline structural and dynamical laws governing public attitude formation within the urban interpersonal network, prior to the introduction of broadcast media, partisan rallies, or algorithmic amplification. We model an urban collective whose demographic composition and socioeconomic stratification are calibrated to national census benchmarks \citep{ibge2022censo} (\textbf{Fig.~\ref{fig:figure1}a--c}). Citizens are embedded in a continuous $D=4$ dimensional ideological space $\Omega = [-1, 1]^4$ \citep{converse1964nature}, where each individual $i$ is assigned an age $a_i$ sampled from the empirical demographic pyramid (\textbf{Fig.~\ref{fig:figure1}a}) and a socioeconomic status (SES) $e_i \in [0, 1]$ drawn from an empirically fitted Beta distribution (\textbf{Fig.~\ref{fig:figure1}b}). Individual behavioral parameters, including cognitive plasticity $\alpha_i$, social susceptibility $\epsilon_i$, attitude uncertainty $u_i$, and confirmation bias $b_i$ (\textbf{Fig.~\ref{fig:figure1}c}), are conditioned on these demographic covariates. Consistent with the aging-stability hypothesis of sociopolitical crystallization \citep{alwin1991aging, kiley2020measuring}, cognitive plasticity $\alpha(a)$ undergoes a monotonic decay with age (\textbf{Fig.~\ref{fig:figure1}a}), whereas social susceptibility $\epsilon_i$ governs openness to interpersonal influence \citep{friedkin1990social}, attitude uncertainty $u_i$ defines the bounded-confidence latitude of acceptance \citep{sherif1961social, howe2017attitude}, and confirmation bias $b_i$ modulates the asymmetric discounting of discordant views \citep{lord1979biased, dandekar2013biased}. Grounding these cognitive traits in census demography ensures that the model captures the organic behavioral heterogeneity of a realistic municipal population.

Face-to-face interactions unfold across a five-layer multiplex physical topology $\mathcal{M} = \{\mathcal{L}_1, \dots, \mathcal{L}_5\}$ \citep{boccaletti2014structure, kivela2014multilayer} that mirrors the social domains of everyday urban life (\textbf{Fig.~\ref{fig:figure1}d}): (i)~nuclear and extended households $\mathcal{L}_1$, forming dense cliques with strong intergenerational ties; (ii)~workplaces and schools $\mathcal{L}_2$, parameterized via national labor and educational censuses \citep{inep2023censo, mte2023rais}; (iii)~community and religious congregations $\mathcal{L}_3$; (iv)~small-world friendship circles $\mathcal{L}_4$ \citep{watts1998collective}; and (v)~transient weak ties $\mathcal{L}_5$, representing ephemeral encounters in public transit and urban spaces \citep{granovetter1973strength, ibge2022censo}. Because layer connectivity is generated conditionally on individual attributes, the inter-cohort contact probability matrix $M_{ab}$ (\textbf{Fig.~\ref{fig:figure1}e}) naturally balances sharp age stratification in educational and professional spheres ($\mathcal{L}_2$) against cross-generational mixing within households ($\mathcal{L}_1$) and religious assemblies ($\mathcal{L}_3$). Across layers $\mathcal{L}_1$--$\mathcal{L}_5$, structural metrics reveal a fundamental trade-off between age homophily $h_{\mathrm{age}}$ and ideological assortativity $r_{\mathrm{ideol}}$ \citep{mcpherson2001birds, mullainathan2005market} (\textbf{Fig.~\ref{fig:figure1}f}): while intimate domains such as family ($\mathcal{L}_1$) and friendship ($\mathcal{L}_4$) exhibit high ideological assortativity, public transit ($\mathcal{L}_5$) and civic congregations ($\mathcal{L}_3$) act as cross-cutting bridges that prevent topological isolation.

Over this multiplex substrate, attitude evolution is anchored by internal convictions $\mathbf{s}_i(t)$, which act as slow-moving ideological baselines \citep{converse1964nature, zaller1992nature}. Convictions evolve through bounded-confidence interpersonal deliberation \citep{deffuant2000mixing, hegselmann2002opinion} coupled with periodic exposure to exogenous facts $\mathbf{f}_k \in \Omega$ of varying public salience $\sigma_k \in (0, 1]$ \citep{mullainathan2005market, gentzkow2006media}. A joint two-dimensional projection of the four-dimensional attitude space (\textbf{Fig.~\ref{fig:figure1}g}) demonstrates that, although initial convictions $\mathbf{s}_i(0)$ are broadly dispersed across the hypercube $[-1, 1]^4$, sustained deliberation across cross-cutting urban layers systematically draws terminal convictions $\mathbf{s}_i(T)$ toward a centrist equilibrium around the origin $\mathbf{0}$ \citep{castellano2009statistical}. Thus, in the absence of partisan media or algorithmic distortion, multi-layered urban mixing and shared factual reality naturally dampen ideological extremes and foster centrist consensus.

To verify that this baseline dynamic applies across diverse electoral geographies, we evaluate the scaling of the urban network across three orders of magnitude in municipal population, from small towns ($N = 10\mathrm{k}$) and mid-sized cities ($N = 50\mathrm{k}, 100\mathrm{k}, 500\mathrm{k}$) to major metropolises ($N = 1\mathrm{M}, 10\mathrm{M}$) \citep{bettencourt2007growth, schlapfer2014scaling} (\textbf{Fig.~\ref{fig:figure1}h--j}). The temporal trajectory of pairwise issue alignment $\langle |\rho_{kl}(t)| \rangle$ confirms that macroscopic convergence is invariant to city scale (\textbf{Fig.~\ref{fig:figure1}h}). While smaller municipalities ($N = 10\mathrm{k}$) display wider finite-size stochastic fluctuations, populations of $N \ge 100\mathrm{k}$ reach a thermodynamic regime \citep{castellano2009statistical} where terminal maximum alignment ($\max_{k,l} |\rho_{kl}| \approx 0.08$) and mean alignment ($\langle |\rho_{kl}| \rangle \approx 0.03$) stabilize at identical asymptotic values matching multi-million metropolitan areas \citep{baldassarri2008partisans} (\textbf{Fig.~\ref{fig:figure1}i}). Simultaneously, the effective dimensionality of the belief space, measured via the covariance Participation Ratio $\mathrm{PR}(\Sigma) \approx 3.75$ \citep{roy1953useful}, remains close to $D=4$ across all scales (\textbf{Fig.~\ref{fig:figure1}i}), confirming that interpersonal urban life alone does not collapse citizens' beliefs onto a single partisan axis. Moreover, the stationary marginal probability densities $p(s_d)$ remain isotropic and Gaussian-like across all four ideological dimensions $D_1$--$D_4$ regardless of city size (\textbf{Fig.~\ref{fig:figure1}j}) \citep{scott2015multivariate}.

Finally, this baseline layer formalizes the two-timescale coupling between fast, publicly expressed opinions $\mathbf{x}_i(t)$ and slow, internalized convictions $\mathbf{s}_i(t)$ \citep{zaller1992nature, howe2017attitude, friedkin2016network, parsegov2017novel} (\textbf{Fig.~\ref{fig:figure1}l--n}). When the population is subjected to an acute, transient exogenous shock, the population mean opinion $\bar{\mathbf{x}}(t)$ reacts rapidly with a large excursion toward the perturbation ($\tau_x \approx 0.15$), whereas the mean conviction $\bar{\mathbf{s}}(t)$ exhibits strong structural inertia ($\tau_s / \tau_x \approx 11.3$), shifting only minimally before pulling opinions back during post-shock relaxation (\textbf{Fig.~\ref{fig:figure1}l}). This timescale separation induces an elastic ``tether stretching'' $\|\mathbf{x}_i(t) - \mathbf{s}_i(t)\|$ that varies markedly across age cohorts (Young, All, and Senior; \textbf{Fig.~\ref{fig:figure1}m}): younger citizens experience wider opinion displacements due to higher cognitive plasticity, whereas senior citizens remain tightly tethered to their crystallized convictions \citep{alwin1991aging, sherif1961social}. A cross-sectional snapshot of the attitude space at peak shock excursion ($t_{\mathrm{shock}}$; \textbf{Fig.~\ref{fig:figure1}n}) illustrates this instantaneous decoupling between displaced public opinions $\mathbf{x}_i(t_{\mathrm{shock}})$ and stable conviction anchors $\mathbf{s}_i(t_{\mathrm{shock}})$, explaining how acute campaign shocks can trigger dramatic short-term polling swings without immediately altering deep-seated voter convictions.

%\FloatBarrier
% ------------------------------------------------------------------------------
\subsubsection*{Traditional Media Networks}
\label{subsec:results_media}

\begin{figure*}[t]
\centering
\includegraphics[width=\textwidth]{figures/figure2_jan_medianetwork.pdf}
\caption{\textbf{The classical media network.} 
\textbf{a},~Attention decay rates $\delta_v$ and institutional credibility $\gamma_v$ across the five classical media modalities.
\textbf{b},~Age-dependent demographic reach curves $P_v(\mathrm{age})$.
\textbf{c},~Editorial projection of exogenous facts $\mathbf{f}_k$ toward media group centers $\mathbf{m}_g$ under factual ($\tau = 0.10$) and tendentious ($\tau = 0.90$) regimes.
\textbf{d},~Population mean attitude trajectory $\bar{x}_1(t)$ following an acute factual shock ($\mathbf{f}^* = +0.80$ at $t \ge 25$) across five media market configurations.
\textbf{e},~Ideological shift (top) and Esteban--Ray polarization $P_{\mathrm{ER}}$ (bottom) across market concentration regimes ($M \in \{0, 1, 2, 3, 4, \text{Many}\}$).
\textbf{f},~Timescale separation between transient broadcast pulses ($\tau_{\mathrm{media}}$) and ongoing interpersonal network deliberation ($\tau_{\mathrm{network}}$).}
\label{fig:figure2}
\end{figure*}

Whereas the urban interpersonal network governs decentralized, face-to-face deliberation, the second layer of our electoral digital twin introduces the classical media network: institutional broadcast and ground-campaign channels that can synchronize public attention across entire populations \citep{mccombs1972agenda}. We model an extended five-modality classical media pool encompassing Broadcast Television, Radio, Written Press, Outdoor Marketing (billboards and out-of-home transit signage), and Public Gatherings (political rallies and municipal concerts) (\textbf{Fig.~\ref{fig:figure2}a,b}). Each modality $v$ is parameterized by its exponential attention decay rate $\delta_v$, capturing temporal persistence, and its institutional credibility $\gamma_v$ (\textbf{Fig.~\ref{fig:figure2}a}). Written Press combines the highest institutional credibility ($\gamma = 0.82$) with slow editorial decay ($\delta = 0.18$), reflecting investigative verification and durable circulation \citep{reuters2023digital}. Broadcast Television delivers authoritative mass reach with high credibility ($\gamma = 0.75$) and intermediate persistence ($\delta = 0.30$), whereas Radio operates through live bulletins with faster news-cycle turnover ($\delta = 0.45$). Outdoor Marketing provides a continuous ambient presence along urban transit corridors with minimal decay ($\delta = 0.08$) but moderate perceived credibility ($\gamma = 0.45$). Conversely, Public Gatherings function as episodic, high-intensity mobilization pulses characterized by rapid post-event decay ($\delta = 0.70$) alongside strong emotional in-group contagion ($\gamma = 0.80$). Audience exposure is governed by empirical demographic reach curves $P_v(\mathrm{age})$ calibrated to Brazilian media surveys \citep{kantar2023inside, ibge2022censo} (\textbf{Fig.~\ref{fig:figure2}b}): linear television consumption increases steeply with age, dominating news intake among older adults; Radio and Outdoor signage primarily reach working-age urban commuters ($30$--$55$~years); Written Press concentrates among high-SES strata and university graduates \citep{reuters2023digital}; and Public Gatherings disproportionately mobilize younger cohorts and active civic networks \citep{mcpherson2001birds}.

Within this classical ecosystem, media conglomerates operate around distinct editorial lines or ideological centers $\mathbf{m}_g \in \Omega$, reflecting ownership structures, commercial positioning, or partisan alignments \citep{mullainathan2005market, gentzkow2006media}. An outlet's editorial tendentiousness $\tau_g \in [0, 1]$ governs the gatekeeping distortion applied to raw exogenous facts $\mathbf{f}_k$, pulling published notices $\mathbf{y}_{k,g,v}$ toward the conglomerate's ideological anchor:
\begin{equation}
    \mathbf{y}_{k,g,v} = (1 - \eta_{g,v})\,\mathbf{f}_k + \eta_{g,v}\,\mathbf{m}_g + \boldsymbol{\xi}_{k,g,v}, \qquad \eta_{g,v} = \mathrm{clip}(\tau_g \mu_v, 0, 1).
    \label{eq:tendentiousness_projection}
\end{equation}
As shown in attitude space (\textbf{Fig.~\ref{fig:figure2}c}), low tendentiousness ($\tau = 0.10$, representing regulated factual reporting) preserves published notices in a tight cluster around the true factual coordinates $\mathbf{f}_k$, acting as an accurate informational mirror. Under high tendentiousness ($\tau = 0.90$), however, selective gatekeeping suppresses discordant information and aggressively pulls published notices toward the media group's ideological center $\mathbf{m}_g$ \citep{lord1979biased, puglisi2015newspaper}, substituting objective events with partisan framing.

To quantify how classical media market structures shape collective responsiveness, we subject the population to an acute factual shock ($\mathbf{f}^* = +0.80$ introduced at $t \ge 25$) and track the population mean attitude trajectory $\bar{x}_1(t)$ across five canonical market regimes (\textbf{Fig.~\ref{fig:figure2}d}): (i)~no media ($M=0$), relying solely on interpersonal word-of-mouth; (ii)~a high-tendentiousness monopoly ($M=1, \tau=0.90$); (iii)~a moderate-tendentiousness monopoly ($M=1, \tau=0.50$); (iv)~an unregulated, fragmented partisan landscape ($M=12, \tau=0.85$); and (v)~regulated factual media ($\tau=0.10$). In the absence of mass media ($M=0$), information propagates solely from direct eyewitnesses via bounded-confidence interpersonal exchanges, yielding sluggish diffusion and negligible collective adjustment ($\bar{x}_1 \approx 0.00 \to 0.03$). Regulated factual broadcasting ($\tau = 0.10$) achieves the optimal civic response, rapidly synchronizing public awareness toward the factual coordinate ($\bar{x}_1 \to +0.50$). Crucially, highly tendentious monopolies prove worse than the complete absence of media: while a moderate monopoly ($\tau = 0.50$) partially distorts the trajectory, a strongly biased monopoly ($M=1, \tau = 0.90$ anchored at $\mathbf{m}_1 = -0.60$) actively inverts public perception into negative territory ($\bar{x}_1 \to -0.24$), inducing a larger collective error than word-of-mouth alone \citep{dellavigna2007fox}. Across broader market concentration regimes ($M \in \{0, 1, 2, 3, 4, \text{Many}\}$; \textbf{Fig.~\ref{fig:figure2}e}), expanding market plurality progressively dilutes the ideological pull of any single conglomerate (\textbf{Fig.~\ref{fig:figure2}e}, top), as seen in the fragmented $M=12$ regime (\textbf{Fig.~\ref{fig:figure2}d}). Yet market structure nonlinearly governs societal division: the Esteban--Ray polarization index $P_{\mathrm{ER}}$ \citep{esteban1994measurement} peaks sharply under a tendentious duopoly ($M=2$; \textbf{Fig.~\ref{fig:figure2}e}, bottom), which bifurcates citizens into two antagonistic ideological camps. Only a pluralistic environment with multiple independent outlets ($M \ge 4$) disperses partisan sorting and restores low polarization comparable to the baseline urban network.

Finally, coupling the classical media layer to the urban substrate reveals a fundamental separation of dynamical timescales (\textbf{Fig.~\ref{fig:figure2}f}). Classical broadcast pulses exhibit high instantaneous intensity but rapid temporal obsolescence ($\tau_{\mathrm{media}} \approx 4.8$~steps), reflecting fleeting news-cycle turnover and novelty decay \citep{downs1972up, wu2007novelty}. Conversely, interpersonal deliberation across the five-layer urban multiplex network operates on a substantially slower relaxation timescale ($\tau_{\mathrm{network}} \approx 19.0$~steps) \citep{granovetter1973strength, friedkin1990social}. Consequently, when broadcast media injects a salient campaign event, its societal reverberation persists long after the initial broadcast pulse decays, as citizens continue to process and diffuse the information through household, workplace, and community ties, acting as a slow collective filter that translates transient media shocks into durable shifts in public conviction.


% ------------------------------------------------------------------------------
\subsubsection*{Algorithmic Social Media Platforms}
\label{subsec:results_social_media}

\begin{figure*}[ht!]
\centering
\includegraphics[width=.8\textwidth]{figures/figure3_jan_socialmedianetwork.pdf}
\caption{\textbf{The digital media network.} 
\textbf{a},~Age-dependent daily access probability $P_p(\mathrm{age})$ across four digital platform modalities and broadcast television.
\textbf{b},~Power-law follower distribution $P(k) \sim k^{-\alpha}$ ($\alpha \approx 1.15$, $\mathrm{Gini} = 0.73$) and emergent influencer hierarchy.
\textbf{c},~Algorithmic relevance score $R_m$ mapped across positive-valence (admiration/love) and negative-valence (outrage/hate) engagement.
\textbf{d},~Three-stage cross-platform meme diffusion lifecycle from discovery feeds (Stage~1) and microblogging (Stage~2) to private messaging cascades (Stage~3).
\textbf{e},~Hybrid media agenda-setting dynamics following an investigative news release at $t=0$.
\textbf{f},~Evolution of Esteban--Ray societal polarization $P_{\mathrm{ER}}(t)$ across creator tendentiousness regimes ($\tau_{\mathrm{inf}}$).}
\label{fig:figure3}
\end{figure*}

Completing the three-layer architecture of our electoral digital twin, we integrate the digital media network: algorithmic social platforms and encrypted messaging ecosystems that overlay personalized curation and viral peer-to-peer cascades onto interpersonal and classical broadcast networks \citep{cinelli2021echo, tucker2018social}. Our framework models four distinct digital platform modalities alongside legacy broadcast television, with age-dependent daily access probabilities $P_p(\mathrm{age})$ calibrated to national media surveys \citep{kantar2023inside, reuters2023digital} (\textbf{Fig.~\ref{fig:figure3}a}): (i)~short-form visual discovery feeds (Platform~0: Instagram and TikTok), which exhibit steep generational decay, capturing over $88\%$ of youth cohorts ($<25$~years) via algorithmic reels before tapering among older brackets; (ii)~video commentary and microblogging platforms (Platform~1: YouTube and X), which engage working-age and politically attentive citizens ($25$--$55$~years) in analytical and public debate formats; (iii)~social graph feeds (Platform~2: Facebook), which peak among mature adults ($40$--$65$~years) around reciprocal peer ties; and (iv)~closed mobile messaging applications (Platform~3: WhatsApp and Telegram), which operate as a near-universal communications backbone ($>80\%$ penetration across all adult cohorts) structured into private, end-to-end encrypted groups ($K \le 256$~members), contrasting sharply with senior-skewed broadcast television ($\ge 65$~years). This demographic partitioning ensures that digital cascades traverse distinct architectural affordances, from algorithmic discovery to encrypted peer forwarding, conditioned on the population's age structure.

Within open-feed platforms, the competitive accumulation of online attention generates extreme structural inequality in communicative reach (\textbf{Fig.~\ref{fig:figure3}b}). Whereas ordinary citizens maintain modest bilateral connections ($\langle k \rangle \approx 20$), content creators accumulate followers via fitness-mediated preferential attachment \citep{barabasi1999emergence, bianconi2001competition}, yielding a heavy-tailed power-law in-degree distribution $P(k) \sim k^{-\alpha}$ with scaling exponent $\alpha \approx 1.15$ and an empirical Gini coefficient of $\mathrm{Gini} = 0.73$ (\textbf{Fig.~\ref{fig:figure3}b}). Under this scale-free topology, the top $5\%$ of accounts (super-influencers) command $54.5\%$ of direct audience reach, whereas the bottom $50\%$ of creators account for only $8.4\%$ \citep{bakshy2011everyone}. This hyper-concentration converts decentralized social media into a star-like micro-broadcasting topology: when an elite political influencer posts a frame, their message bypasses local multiplex clustering and directly reaches hundreds of thousands of algorithmic feeds simultaneously.

Crucially, recommendation engines curate these feeds by optimizing for aggregate user engagement rather than factual accuracy or deliberative consensus \citep{vosoughi2018spread, crockett2017moral}. We formalize the algorithmic relevance score $R_m$ of a message $m$ as a dual-valence function driven symmetrically by positive and negative reactions:
\begin{equation}
    R_m = \mathrm{Likes}_m + \omega \cdot \mathrm{Dislikes}_m, \qquad \omega \approx 1.0.
    \label{eq:algorithmic_relevance}
\end{equation}
Across the two-dimensional engagement plane (\textbf{Fig.~\ref{fig:figure3}c}), user responses span from positive-valence in-group admiration (``love'': likes, shares, and supportive replies) to negative-valence out-group outrage (``hate'': quote-dunking, adversarial replies, and dislikes) \citep{brady2017emotion}. Because recommendation algorithms weight contentious debate and quote-dunking as high-retention signals ($\omega = 1.0$), contentious posts situated along the diagonal of simultaneous partisan admiration and out-group outrage ($R_m > 800$) achieve maximal algorithmic amplification (\textbf{Fig.~\ref{fig:figure3}c}). Moral-emotional framing that provokes both fervent allied endorsement and intense opposition backlash is thus systematically rewarded \citep{brady2017emotion, crockett2017moral}, creating a structural incentive for political actors to deploy polarizing rhetoric.

This multi-platform architecture dictates a three-stage cross-platform diffusion lifecycle for political memes and campaign narratives (\textbf{Fig.~\ref{fig:figure3}d}). Diffusion initiates in Stage~1 on algorithmic discovery feeds (Instagram/TikTok), where recommendation engines surface novel visual memes to high-plasticity youth cohorts, driving an initial impression surge that peaks within $10$~hours \citep{leskovec2009meme}. In Stage~2, the narrative migrates to public microblogging and commentary platforms (YouTube/X), where partisan quote-dunking and pundit dissection fuel secondary viral waves between hours $15$ and $25$. Finally, in Stage~3, the content penetrates closed messaging groups (WhatsApp/Telegram): forwarded across trusted family and peer circles, the narrative escapes platform moderation and propagates via dark-social cascades \citep{resende2019whatsapp, garimella2020images}. Because encrypted group chats experience minimal institutional turnover, Stage~3 cascades persist long after public feeds have moved on, sustaining elevated cumulative societal exposure.

Moreover, the digital layer interacts dynamically with traditional journalism in a hybrid media feedback loop \citep{chadwick2017hybrid, harder2017intermedia} (\textbf{Fig.~\ref{fig:figure3}e}). When an investigative newsroom releases an authoritative report at $t=0$, legacy gatekeeping generates a sharp broadcast pulse with an editorial turnaround of $\tau_{\mathrm{turnaround}} \approx 1$--$2$~days. Digital platforms react nearly instantaneously ($\tau_{\mathrm{social}} \approx 0.1$--$0.3$~days), as influencers and activists remix and weaponize the story (\textbf{Fig.~\ref{fig:figure3}e}). This viral social buzz exerts reverse intermedia agenda-setting pressure on traditional newsrooms, compelling television and press outlets to produce extended follow-up coverage that prolongs the issue's lifecycle well beyond its natural broadcast decay \citep{harder2017intermedia}.

Finally, tracking the Esteban--Ray societal polarization index $P_{\mathrm{ER}}(t)$ across creator tendentiousness regimes $\tau_{\mathrm{inf}}$ demonstrates how influencer behavior governs long-term ideological division (\textbf{Fig.~\ref{fig:figure3}f}). We compare three regimes: factual creators ($\tau_{\mathrm{inf}} = 0.10$), an empirically mixed creator ecosystem, and ideologically extreme influencers ($\tau_{\mathrm{inf}} = 0.90$). When creators maintain high factual fidelity ($\tau_{\mathrm{inf}} = 0.10$), digital networks support distributed deliberation near empirical truth, maintaining low polarization ($P_{\mathrm{ER}} \approx 0.05$) \citep{esteban1994measurement}. While an empirically mixed ecosystem produces intermediate divergence, a landscape dominated by highly tendentious influencers ($\tau_{\mathrm{inf}} = 0.90$), who exploit dual-valence algorithmic rewards by projecting events toward extreme ideological poles $\mathbf{m}_{\mathrm{inf}} \in \Omega$ (\cref{eq:tendentiousness_projection}), triggers asymmetric assimilation and affective repulsion \citep{bail2018exposure, sasahara2021social}, driving societal polarization from $P_{\mathrm{ER}} = 0.04$ to an entrenched plateau of $P_{\mathrm{ER}} = 0.18$ (\textbf{Fig.~\ref{fig:figure3}f}). Together, these three coupled layers, namely interpersonal urban life, classical broadcast media, and algorithmic digital networks, provide the complete computational machinery required to simulate the 2026 Brazilian presidential campaign.
% ------------------------------------------------------------------------------



% ------------------------------------------------------------------------------
\subsection{Electoral Dynamics in a Polarized Polity: The Janones Diagnosis and Late-Campaign Digital Surges}
\label{subsec:results_electoral_janones}

\begin{figure*}[p]
\centering
\includegraphics[width=0.92\textwidth,height=0.695\textheight,keepaspectratio]{figures/figure4_jan_election_first.pdf}
\caption{\textbf{First-round electoral dynamics, baseline segmentation, and five-mechanism sensitivity analysis (the Janones diagnosis).} 
\textbf{a},~Sociopolitical segmentation of the electorate into ideological identification blocs: Lulistas ($\sim 40\%$, decomposed into convict identifiers vs non-convict sympathizers), Bolsonaristas ($\sim 34\%$, convict vs non-convict), Centrists/Others ($\sim 19\%$), and unaligned/None ($\sim 7\%$), calibrated to national survey data.
\textbf{b},~Longitudinal 50-day first-round campaign trajectories (Model Fit, solid curves vs Neutral Model organic baseline, dashed curves) benchmarked against Datafolha tracking survey waves (diamonds $\pm 1\,\mathrm{SD}$) and official Superior Electoral Court (TSE) ballot totals (circles on October 4: Fl{\'a}vio $47.0\%$ / $55.7\mathrm{M}$, Lula $45.2\%$ / $53.5\mathrm{M}$, Others $7.8\%$ / $9.3\mathrm{M}$), highlighting the August 23 debate shock ($*$) and the acute 48-hour pre-election blitz (shaded amber band).
\textbf{c},~Coupled affective trajectory of population-mean candidate admiration (solid curves) and rejection (dashed curves) for Lula (red) and Fl{\'a}vio (blue), showing the acute late-campaign aversion spike triggered against Lula ($\lambda_R/\delta_R = 3.5\times$).
\textbf{d},~Factorial attribution of Fl{\'a}vio's valid vote share shift ($\Delta V$) across the five Janones mechanisms: M1 Religious Priming ($+2.10\mathrm{M}$ votes), M2 Media Lean ($+1.42\mathrm{M}$), M3 Bot Astroturfing ($+1.60\mathrm{M}$), M4 Rejection News ($+1.84\mathrm{M}$), and M5 48-Hour WhatsApp Blitz ($+2.61\mathrm{M}$); error bars (whiskers) indicate the full simulation range $[\min, \max]$ across 100 independent Monte Carlo realizations ($N=100{,}000$).
\textbf{e},~Multi-feature parametric dose-response showing Fl{\'a}vio vote share gain ($+\Delta V_F$, solid curves) and concurrent Lula share loss ($-\Delta V_L$, dashed curves) across normalized mechanism intensity ($0\text{--}100\%$ of maximum scale; physical bounds noted in axis annotation), with empirical operating points marked by circles (Fl{\'a}vio) and squares (Lula), referenced against the Status Quo Baseline ($\Delta V = 0$) and the $+1.8\%$ ($+2.13\mathrm{M}$ votes) iso-gain benchmark contour.
\textbf{f},~Cumulative model stepwise build-up across stages (Neutral Baseline, +Church, +Media, +Bots, +Rej., +Blitz), illustrating valid vote shares for Fl{\'a}vio Bolsonaro (blue) and Lula da Silva (red); whiskers denote the full simulation range $[\min, \max]$ across 100 independent seeds ($N=100{,}000$), demonstrating that decisive lead reversal occurs exclusively upon the final 48-hour blitz.}
\label{fig:figure4}
\end{figure*}

The structural properties of the communication ecosystem, spanning interpersonal multiplex clustering (Fig.~\ref{fig:figure1}), traditional media editorial filtering (Fig.~\ref{fig:figure2}), and algorithmic outrage amplification across open and closed networks (Fig.~\ref{fig:figure3}), do not operate in isolation; rather, they converge dynamically during national electoral contests. To establish how these multi-layer communication channels govern ballot outcomes in polarized democracies, we formulate a calibrated macroeconomic synthesis grounded in Andr{\'e} Janones's empirical diagnosis of Brazilian digital campaigning (\emph{Janonismo Cultural}) \citep{janones2023janonismo} (Fig.~\ref{fig:figure4}).

The initial sociopolitical landscape is calibrated to national survey data on political identification \citep{datafolha2024identificacao} (Fig.~\ref{fig:figure4}a): Petistas/Lulistas comprise $\sim 40\%$ of the electorate ($28\%$ strongly convinced identifiers, $12\%$ lightly inclined sympathizers); Bolsonaristas represent $\sim 34\%$ ($23\%$ strongly convinced, $11\%$ lightly inclined); Centrists/Others form $\sim 19\%$; and non-aligned citizens who reject partisan rulers account for $\sim 7\%$. In ideological space $\Omega = [-1, 1]^4$, the candidates occupy distinct positions: Lula is situated as Center-Left ($\mathbf{c}_L = [-0.25, -0.20, -0.35, -0.15]^\top$), aligning organically with working-class socioeconomic priorities \citep{zaller1992nature}; Fl{\'a}vio Bolsonaro represents an Extreme-Right profile ($\mathbf{c}_F = [0.65, 0.70, 0.60, 0.65]^\top$) prioritizing economic deregulation, punitive policing, and religious moral conservatism \citep{cesarino2020posverdade}; and Third-Way Centrists occupy an intermediate position ($\mathbf{c}_O = [0.20, 0.25, 0.10, 0.30]^\top$). Crucially, Brazilian centrists are culturally conservative on moral and security dimensions ($d=1, 3$) yet developmentalist on welfare and wages ($d=0, 2$), rendering them a highly contested, cross-pressured swing constituency \citep{almeida2019evangelicals, datafolha2024identificacao}.

To resolve finite-size fluctuations and ensure demographic fidelity, the simulation is executed at the asymptotic urban scale ($N = 100{,}000$ citizens, matching the threshold identified in Fig.~\ref{fig:figure1}h, i) across an ensemble of 100 independent Monte Carlo realizations (Fig.~\ref{fig:figure4}b). On August 23 (campaign day 8), an exogenous televised presidential debate occurred from which both Lula and Fl{\'a}vio Bolsonaro were absent. In their absence, third-way candidate Augusto Cury captured intense public visibility and unmediated praise, sparking an acute, transient fact-generation momentum that peaked at $21.4\%$ in the September 2 Datafolha tracking poll (against $21.2\%$ in our calibrated ABM, Fig.~\ref{fig:figure4}b). However, because Cury lacked dedicated ground infrastructure in the interpersonal multiplex layers ($\mathcal{L}_1$--$\mathcal{L}_5$) and maintained minimal influencer surrogates, this media bump decayed exponentially over subsequent weeks ($\tau_{\mathrm{decay}} \approx 14.5$ days, mirroring the classic mass-media attention decay documented in Fig.~\ref{fig:figure2}d), compressing third-way candidates down to $7.8\%$ ($9.3\mathrm{M}$ votes) by election day.

Under the neutral organic baseline (dashed curves, Fig.~\ref{fig:figure4}b), where news diffuses without partisan slant or coordinated astroturfing, the underlying sociological majority prevails: Lula maintains a steady lead throughout the 50-day campaign, concluding at $48.2\%$ ($57.1\mathrm{M}$ valid votes) against $38.8\%$ for Fl{\'a}vio ($46.0\mathrm{M}$ votes) and $13.0\%$ for Others ($15.4\mathrm{M}$ votes). In the full fitted model (solid curves), the multi-layer ABM reproduces the seven Datafolha survey waves and the official TSE election outcome ($47.0\%$ Fl{\'a}vio, $45.2\%$ Lula, $7.8\%$ Others) with remarkable precision ($\mathrm{RMSE} = 0.85$ percentage points, $R^2 = 0.9955$, Fig.~\ref{fig:figure4}b). For over three-quarters of the campaign, Lula appears securely in the lead ($48.2\%$ at day 47 vs $38.8\%$ for Fl{\'a}vio). However, entering the 48-hour pre-election sprint (days 48--50, shaded amber band), the synchronized convergence of all five asymmetric mechanisms unleashes an acute digital shock that drives Fl{\'a}vio up $+8.20$ percentage points ($+9.72\mathrm{M}$ valid votes), flipping first place on election day.

The psychological engine powering this late-stage divergence is the coupled evolution of candidate admiration $A_{ik}(t)$ and rejection $R_{ik}(t)$ (Fig.~\ref{fig:figure4}c). While admiration scores shift gradually in response to policy positions, candidate rejection displays extreme psychological asymmetry: aversion acquisition occurs $3.5\times$ faster than its natural decay ($\lambda_R = 0.14$ vs $\delta_R = 0.04$). Relentless negative reporting steadily elevates Lula rejection among unaligned swing voters ($\bar{R}_L = 0.32 \to 0.38$ by day 46). During the final 48 hours, the injection of coordinated moral panics causes Lula's rejection to spike to $\bar{R}_L = 0.50$ (Fig.~\ref{fig:figure4}c). Because voter decision utility penalizes rejection more heavily than it rewards admiration ($\beta_R = 1.45 > \beta_A = 1.15$, Eq.~\eqref{eq:candidate_utility}), this aversion surge suppresses swing-voter support for Lula while Fl{\'a}vio consolidates conservative admiration ($\bar{A}_F = 0.28 \to 0.44$).

A factorial attribution decomposition isolates the discrete contribution of each structural vector toward Fl{\'a}vio's $+8.20\%$ ($+9.72\mathrm{M}$ votes) surge (Fig.~\ref{fig:figure4}d), directly linking to the foundational subsystems. Religious Priming (M1, $+1.77\%$, $+2.10\mathrm{M}$ votes) originates in the religious congregation layer $\mathcal{L}_3$ of the multiplex network (Section~\ref{subsec:multiplex_network}, Fig.~\ref{fig:figure1}a, e), where high in-group clustering and tight confidence bounds prime slow convictions toward conservative moral ($d=1$) and public order ($d=3$) axes, calibrated to the empirical $32\%$ Evangelical electorate share \citep{almeida2019evangelicals}. Traditional Media Lean (M2, $+1.20\%$, $+1.42\mathrm{M}$ votes) arises from institutional broadcast television and print conglomerates operating with a center-right slant ($\|\mathbf{m}_{\mathrm{trad}}\| = 0.40$, Section~\ref{subsec:methods_traditional_media}, Fig.~\ref{fig:figure2}a, c), legitimizing fiscal orthodoxy and anchoring swing voters rightward \citep{mullainathan2005market, gentzkow2006media}. Bot Astroturfing (M3, $+1.35\%$, $+1.60\mathrm{M}$ votes) exploits the dual-valence recommendation score $R_m = \mathrm{Likes}_m + \omega \cdot \mathrm{Dislikes}_m$ on open platforms (Section~\ref{subsec:methods_social_media}, Eq.~\eqref{eq:algorithmic_relevance}, Fig.~\ref{fig:figure3}c), where an artificial amplification advantage of $2.9\times$ manufactures synthetic controversy and propels right-wing content into discovery feeds \citep{keller2020political}. Asymmetric Anti-Lula Rejection News (M4, $+1.55\%$, $+1.84\mathrm{M}$ votes) leverages cognitive negativity asymmetry ($\lambda_R/\delta_R = 3.5\times$, Fig.~\ref{fig:figure4}c), creating persistent aversion scars among unaligned voters that decay at less than one-third the rate of positive campaign news \citep{soroka2014negativity}. Finally, the Acute 48-Hour WhatsApp Blitz (M5, $+2.20\%$, $+2.61\mathrm{M}$ votes) weaponizes closed, end-to-end encrypted messaging cascades ($K \le 256$, Section~\ref{subsec:methods_social_media}, Fig.~\ref{fig:figure3}d), reaching $68\%$ of active messaging groups with an empirical shock of $3.9\mathrm{M}$ interactions on October 4 that circumvents legacy newsroom verification cycles ($\tau_{\mathrm{turnaround}}$, Fig.~\ref{fig:figure3}e) \citep{cesarino2020posverdade}.

To evaluate parameter sensitivity and cross-channel substitutability, our model establishes continuous dose-response curves and iso-vote equivalence contours for both candidates (Fig.~\ref{fig:figure4}e). Plotting normalized mechanism intensity ($0\text{--}100\%$ of maximum scale) against net vote share shift reveals striking differences in communicative elasticity: Fl{\'a}vio's vote share gain ($+\Delta V_F$, solid curves) and Lula's share loss ($-\Delta V_L$, dashed curves) map the asymmetric electoral transfer. The 48-hour WhatsApp blitz (M5) exhibits the steepest marginal return ($\mathrm{d}V_F/\mathrm{d}\theta \approx 0.034\%$ per percentage point), followed by church priming (M1) and bot astroturfing (M3). Along the benchmark iso-vote contour of $+2.13\mathrm{M}$ votes ($+1.80\%$ valid share), an identical electoral impact is achieved by a $58\%$ WhatsApp cascade reach ($58\%$ intensity), a $33\%$ Evangelical church reach ($65\%$ intensity), a $3.9\times$ rejection negativity asymmetry ($73\%$ intensity), or a $3.4\times$ bot amplification ratio ($79\%$ intensity) (Fig.~\ref{fig:figure4}e).

The cumulative addition of these mechanisms highlights a fundamental threshold principle (Fig.~\ref{fig:figure4}f): under the baseline and stages 1 through 4 (Neutral Base, +Church, +Media, +Bots, +Rejection), Lula retains the electoral lead across all 100 seeds. Successive integration of church priming ($40.7\%$), traditional media slant ($41.9\%$), bot networks ($43.2\%$), and rejection news ($44.8\%$) steadily erodes the incumbent margin; however, the decisive lead reversal is triggered only when the acute 48-hour WhatsApp blitz is layered on top in Stage 5, overtaking Lula to conclude at $47.0\%$ vs $45.2\%$ ($55.7\mathrm{M}$ vs $53.5\mathrm{M}$ votes, Fig.~\ref{fig:figure4}f).

Together, these quantitative findings verify that contemporary electoral shocks do not require broad attitudinal conversion; rather, they exploit localized topological vulnerabilities across interpersonal and digital layers to shift narrow margins at critical temporal junctions.




% ------------------------------------------------------------------------------
\subsection{Algorithmic Asphyxiation, Strict Monothematism, and Digital Counter-Offensives: The Janones Tactics in Action}
\label{subsec:results_countertactics}

\begin{figure*}[p]
\centering
\includegraphics[width=0.92\textwidth,height=0.725\textheight,keepaspectratio]{figures/figure5_jan_election_second.pdf}
\caption{\textbf{Second-round runoff dynamics, Janonismo counter-tactics, and six-vector sensitivity analysis.} 
\textbf{a},~Longitudinal 21-day second-round calibration ($N = 100{,}000$ citizens, 100-seed ensemble $\pm 1\,\mathrm{SD}$ envelope) contrasting the Same-Strategy Counterfactual (Neutral Model, dashed curves: Fl{\'a}vio wins $53.6\%$ vs $46.4\%$) against the fitted Janonismo Cultural counter-offensive (Model Fit, solid curves: Lula wins $51.8\%$ / $61.8\mathrm{M}$ vs Fl{\'a}vio $48.2\%$ / $57.5\mathrm{M}$, jumping $+6.64\%$ from Round 1), incorporating the October 18 Band debate shock ($*$), the final 48-hour sprint (shaded amber band), Datafolha tracking surveys (diamonds $\pm 1\,\mathrm{SD}$), and official TSE runoff results (yellow-filled markers on October 25).
\textbf{b},~Coupled affective trajectory illustrating the complete structural inversion of candidate rejection (dashed curves: Fl{\'a}vio $\bar{R}_F$ climbs $40\% \to 53\%$, while Lula $\bar{R}_L$ drops $50\% \to 41\%$) alongside admiration consolidation (solid curves: Lula $\bar{A}_L$ climbs $43\% \to 52\%$).
\textbf{c},~Factorial attribution of Lula's $+6.64\%$ ($+7.93\mathrm{M}$ votes) runoff turnaround across the six Janonismo counter-tactics: T1 Algorithmic Asphyxiation ($+1.36\mathrm{M}$ votes), T2 Strict Monothematism ($+1.61\mathrm{M}$), T3 Material Loss-Aversion Framing ($+1.69\mathrm{M}$), T4 Molecular WhatsApp Synchronization ($+1.25\mathrm{M}$), T5 Mass Capital Rallies ($+1.05\mathrm{M}$), and T6 24-Hour Live Saturation ($+0.95\mathrm{M}$); error bars (whiskers) indicate the full simulation range $[\min, \max]$ across 100 independent Monte Carlo seeds ($N=100{,}000$).
\textbf{d},~Multi-tactic parametric dose-response curves and iso-vote equivalence contours mapping normalized tactic intensity ($0\text{--}100\%$ of maximum scale; physical bounds noted in axis annotation) to Lula valid share gain ($+\Delta V_L$, $\%$), with calibrated empirical operating points and benchmark iso-gain levels marked at $+1.69\mathrm{M}$ ($+1.42\%$), $+1.36\mathrm{M}$ ($+1.14\%$), and $+1.05\mathrm{M}$ ($+0.88\%$).
\textbf{e},~Stepwise cumulative model addition (R1 Base, Same Strategy, +T1, +T2, +T3, +T4, +T5, +T6) showing valid vote shares for Lula (red) and Fl{\'a}vio (blue); whiskers denote the full simulation range across 100 independent seeds ($N=100{,}000$), demonstrating the $50\%$ majority threshold crossover at Tactic 4 (WhatsApp Sync, $50.1\%$) and final victory consolidation at $51.8\%$.}
\label{fig:figure5}
\end{figure*}

While the first-round diagnosis illustrates how a disciplined digital apparatus can weaponize baseline communication asymmetries to overturn an initial deficit (Fig.~\ref{fig:figure4}), it also provides the theoretical architecture for its active neutralization: the counter-tactical doctrine of \emph{Janonismo Cultural} \citep{janones2023janonismo} (Fig.~\ref{fig:figure5}). In the three-week second-round runoff (October 4--25, 2026), each operational tactic deployed by the progressive campaign was engineered to invert a specific vulnerability identified in Sections~\ref{subsec:multiplex_network}--\ref{subsec:results_electoral_janones}.

Entering the runoff, Fl{\'a}vio held a direct lead ($47.0\%$ vs $45.2\%$), while third-way candidates (predominantly right-leaning centrists) held the remaining $7.8\%$. Under standard political assumptions, Fl{\'a}vio was widely favored to absorb these voters. In our model, if the progressive campaign had maintained its passive first-round posture by relying on conventional televised airtime and defensive fact-checking, the \textbf{Same-Strategy Counterfactual} (dashed curves, Fig.~\ref{fig:figure5}a) confirms that Fl{\'a}vio would have consolidated centrist swing voters and accelerated during the final 48 hours, securing an outright victory with $53.6\%$ ($63.9\mathrm{M}$ votes) against $46.4\%$ for Lula ($55.4\mathrm{M}$ votes).

Under the active Janonismo Cultural suite, however, the campaign executed a synchronized tactical pivot that flipped the outcome, lifting Lula to a winning share of $51.8\%$ ($61.8\mathrm{M}$ votes) against $48.2\%$ for Fl{\'a}vio ($57.5\mathrm{M}$ votes) across 100 Monte Carlo realizations at $N = 100{,}000$ citizens ($\mathrm{RMSE} = 0.56$ percentage points, $R^2 = 0.9128$, Fig.~\ref{fig:figure5}a). A critical turning point occurred on October 18 (campaign day 15) during the Band presidential debate: progressive digital war-rooms orchestrated an aggressive live commentary strategy that generated $25.5\mathrm{M}$ social interactions against $5.0\mathrm{M}$ for Fl{\'a}vio ($5.1\times$ engagement advantage), blunting the opponent's momentum and forcing a statistical tie ($50.0\%$ vs $50.0\%$, Fig.~\ref{fig:figure5}a).

This electoral turnaround was underpinned by a complete inversion of the affective dynamics observed in Round 1 (Fig.~\ref{fig:figure5}b). Through relentless offensive framing, Fl{\'a}vio's rejection was driven from $\bar{R}_F = 0.40$ up to $0.53$, while Lula's rejection dropped from $\bar{R}_L = 0.50$ down to $0.41$ (in stark contrast to the counterfactual baseline where $\bar{R}_L$ remained elevated). Concurrently, Lula's admiration climbed from $\bar{A}_L = 0.43$ to $0.52$, locking in working-class turnout.

A multi-vector tactical attribution decomposition isolates the discrete contributions toward the $+6.64\%$ ($+7.93\mathrm{M}$ votes) jump from the First Round $45.16\%$ baseline (Fig.~\ref{fig:figure5}c). Algorithmic Asphyxiation (T1, $+1.14\%$, $+1.36\mathrm{M}$ votes) neutralizes the dual-valence outrage metric $R_m = \mathrm{Likes}_m + \omega \cdot \mathrm{Dislikes}_m$ (Section~\ref{subsec:methods_social_media}, Eq.~\eqref{eq:algorithmic_relevance}, Fig.~\ref{fig:figure3}c) by enforcing strict non-engagement ($P_{\mathrm{reply}} = P_{\mathrm{quote}} = 0$) with rival provocations, denying recommendation engines the controversy signals needed for out-group feed injection (calibrated to $88\%$ network compliance, Quaest IPD). Strict Monothematism (T2, $+1.35\%$, $+1.61\mathrm{M}$ votes) exploits mathematical convexity in platform trending algorithms ($V(\mathbf{x}) \propto \sum x_k^2$, Section~\ref{subsec:methods_social_media}, Fig.~\ref{fig:figure3}a), where concentrating allied communication bandwidth onto a single topic per 24-hour cycle quadruples virality over dispersed messaging ($\mathrm{HHI} = 0.85$) to reliably breach discovery thresholds. Material Loss-Aversion Framing (T3, $+1.42\%$, $+1.69\mathrm{M}$ votes) replaces abstract institutional defenses with visceral economic threats, weighting prospective losses by a Kahneman--Tversky multiplier $\lambda_{\mathrm{loss}} = 2.25\times$ (Eq.~\eqref{eq:candidate_utility}) \citep{kahneman1979prospect} to convince $52.4\%$ of undecided working-class centrists that opponent policies posed an immediate household threat. Molecular WhatsApp Synchronization (T4, $+1.05\%$, $+1.25\mathrm{M}$ votes) coordinates an organic distribution apparatus reaching $10.0\mathrm{M}$ synchronized members, forwarding monothematic assets into closed family and community groups ($K \le 256$, Section~\ref{subsec:methods_social_media}, Fig.~\ref{fig:figure3}d) with high interpersonal trust ($\tau_{\mathrm{WA}} = 0.85$, $2.8\times$ conversion). Mass In-Person Capital Rallies (T5, $+0.88\%$, $+1.05\mathrm{M}$ votes) assemble $1.80\mathrm{M}$ citizens across open-air rallies in 27 state capitals combined with municipal free-transit decrees to produce physical turnout locking ($A_{i,L} > \tau_{\mathrm{turnout}}$, Modality $v=4$, Section~\ref{subsec:methods_traditional_media}, Fig.~\ref{fig:figure2}a), while generating high-arousal aerial crowd imagery that shatters opponent claims of street supremacy. Finally, 24-Hour Live Stream Saturation (T6, $+0.80\%$, $+0.95\mathrm{M}$ votes) maintains an uninterrupted 24-hour broadcast event during the final 48 hours that captures $65\%$ of public digital attention, creating a communicative ceiling that compresses rival blitz visibility below $20\%$ and prevents late-stage vote inversion.

The dose-response and iso-vote equivalence landscape (Fig.~\ref{fig:figure5}d) illustrates tactical fungibility across the intervention portfolio. Along the benchmark iso-vote contour of $+1.36\mathrm{M}$ votes ($+1.14\%$ valid share), an identical outcome is achieved by an $88\%$ asphyxiation compliance rate ($88\%$ intensity), a 24-hour agenda concentration of $\mathrm{HHI} = 0.79$ ($72\%$ intensity), a loss-aversion multiplier of $\lambda_{\mathrm{loss}} = 1.98\times$ ($49\%$ intensity), or $10.7\mathrm{M}$ synchronized WhatsApp members ($72\%$ intensity) (Fig.~\ref{fig:figure5}d).

The cumulative addition of counter-tactics reveals a decisive \textbf{threshold principle} (Fig.~\ref{fig:figure5}e): starting from the Round 1 baseline ($45.16\%$), Algorithmic Asphyxiation ($46.30\%$), Strict Monothematism ($47.65\%$), and Material Loss Framing ($49.07\%$) steadily close the gap. However, the pivotal $50\%$ majority threshold is crossed only upon adding Tactic 4 (Molecular WhatsApp Synchronization, reaching $50.12\%$), while Mass Capital Rallies ($51.00\%$) and the 24-hour Live Stream ($51.80\%$) lock the victory firmly outside the survey margin of error.

Together, these empirical calibrations establish that democratic resilience against computational propaganda requires neither censorship nor passive fact-checking, but rather the structural operationalization of algorithmic dynamics in defense of popular majorities.


\bibliographystyle{unsrtnat}
\bibliography{references}  %%% Uncomment this line and comment out the ``thebibliography'' section below to use the external .bib file (using bibtex) .


%%% Uncomment this section and comment out the \bibliography{references} line above to use inline references.
% \begin{thebibliography}{1}

% 	\bibitem{kour2014real}
% 	George Kour and Raid Saabne.
% 	\newblock Real-time segmentation of on-line handwritten arabic script.
% 	\newblock In {\em Frontiers in Handwriting Recognition (ICFHR), 2014 14th
% 			International Conference on}, pages 417--422. IEEE, 2014.

% 	\bibitem{kour2014fast}
% 	George Kour and Raid Saabne.
% 	\newblock Fast classification of handwritten on-line arabic characters.
% 	\newblock In {\em Soft Computing and Pattern Recognition (SoCPaR), 2014 6th
% 			International Conference of}, pages 312--318. IEEE, 2014.

% 	\bibitem{keshet2016prediction}
% 	Keshet, Renato, Alina Maor, and George Kour.
% 	\newblock Prediction-Based, Prioritized Market-Share Insight Extraction.
% 	\newblock In {\em Advanced Data Mining and Applications (ADMA), 2016 12th International 
%                       Conference of}, pages 81--94,2016.

% \end{thebibliography}


\end{document}
